\documentclass[11pt,a4paper]{article}
\usepackage[a4paper,margin=27mm,headheight=14pt]{geometry}
\usepackage{fontspec}
\setmainfont{Latin Modern Roman}
\setsansfont{Latin Modern Sans}
\setmonofont{Latin Modern Mono}
\usepackage{amsmath,amssymb,amsthm,mathtools,mathrsfs,enumitem,booktabs,array,xcolor,fancyhdr,longtable}
\definecolor{resultlink}{HTML}{1F4E79}
\definecolor{sectionlink}{HTML}{00656D}
\definecolor{equationlink}{HTML}{643F7A}
\definecolor{citationlink}{HTML}{2F6B3C}
\definecolor{urllink}{HTML}{843C39}
\usepackage[colorlinks=true,linkcolor=sectionlink,citecolor=citationlink,urlcolor=urllink,bookmarksnumbered=true]{hyperref}
\hypersetup{pdftitle={Domains and constructions of Beilinson--Bloch height pairings},pdfauthor={},pdfsubject={Geometric domains from actual Green constructions and reciprocity}}
\newtheorem{theorem}{Theorem}[section]
\newtheorem{proposition}[theorem]{Proposition}
\newtheorem{lemma}[theorem]{Lemma}
\newtheorem{corollary}[theorem]{Corollary}
\newtheorem{hypothesis}[theorem]{Hypothesis}
\theoremstyle{definition}
\newtheorem{definition}[theorem]{Definition}
\newtheorem{example}[theorem]{Example}
\theoremstyle{remark}\newtheorem{remark}[theorem]{Remark}
\numberwithin{equation}{section}
\newcommand{\Q}{\mathbb Q}\newcommand{\R}{\mathbb R}\newcommand{\C}{\mathbb C}
\newcommand{\Z}{\mathbb Z}
\newcommand{\PP}{\mathbb P}\newcommand{\LL}{\mathbb L}
\newcommand{\CH}{\operatorname{CH}}\newcommand{\an}{\mathrm{an}}
\newcommand{\homol}{\mathrm{hom}}\newcommand{\BC}{\mathrm{BC}}
\newcommand{\D}{d'd''}\newcommand{\ob}{\mathfrak o}
\newcommand{\Spec}{\operatorname{Spec}}\newcommand{\Bl}{\operatorname{Bl}}
\newcommand{\Tr}{\operatorname{Tr}}\newcommand{\ch}{\operatorname{ch}}
\newcommand{\alg}{\mathrm{alg}}\newcommand{\geomhom}{\mathrm{geom\,hom}}
\newcommand{\cK}{\mathcal K}\newcommand{\cH}{\mathcal H}
\newcommand{\rk}{\operatorname{rk}}\newcommand{\im}{\operatorname{im}}
\newcommand{\secref}[1]{\hyperref[#1]{\textcolor{sectionlink}{Section~\ref*{#1}}}}
\newcommand{\subref}[1]{\hyperref[#1]{\textcolor{sectionlink}{\S\ref*{#1}}}}
\newcommand{\thmref}[1]{\hyperref[#1]{\textcolor{resultlink}{Theorem~\ref*{#1}}}}
\newcommand{\propref}[1]{\hyperref[#1]{\textcolor{resultlink}{Proposition~\ref*{#1}}}}
\newcommand{\lemref}[1]{\hyperref[#1]{\textcolor{resultlink}{Lemma~\ref*{#1}}}}
\newcommand{\corref}[1]{\hyperref[#1]{\textcolor{resultlink}{Corollary~\ref*{#1}}}}
\newcommand{\hypref}[1]{\hyperref[#1]{\textcolor{resultlink}{Hypothesis~\ref*{#1}}}}
\let\oldeqref\eqref
\renewcommand{\eqref}[1]{\hyperref[#1]{\textcolor{equationlink}{(\ref*{#1})}}}
\newcommand{\proofloc}[1]{\emph{Proof.} See \subref{#1}.}
\setlist{itemsep=3pt,topsep=5pt,leftmargin=2em}
\setlength{\emergencystretch}{2em}
\setlength{\parskip}{3pt}
\pagestyle{fancy}\fancyhf{}\fancyhead[L]{\small Domains of Beilinson--Bloch height pairings}\fancyhead[R]{}\fancyfoot[C]{\thepage}
\title{Domains and constructions of\newline Beilinson--Bloch height pairings}
\author{}
\date{Research manuscript --- revised 23 November 2026}
\begin{document}
\maketitle
\begin{abstract}
We organize the geometric domains of zero-curvature Green forms and normalized
Beilinson--Bloch height pairings. The actual-form, numerical-descent
and reciprocity constructions, together with Chow covariance, turn decompositions and model
corrections into constructions on the exact support complement of a prescribed
cycle. The resulting domains include products of arbitrary curves and
surfaces, their specified motivic summands, Hilbert schemes of points on
surfaces, and the surface-moduli motives furnished by B\"ulles. Other
applications use low-dimensional diagonal support, Chow niveau, flat
fibrations, and semistable models. We give the precise complementary-degree
conditions for global heights, construct Chow identities for separated
hypersurfaces and cyclic covers, and record the field and cover degrees in
the comparison with classical abelian heights. The geometric proofs are given
in full; the separately supplied analytic and model statements are identified
at each use.
\end{abstract}
\tableofcontents

\section{Introduction}\label{sec:intro}

Let $K/\Q_p$ be finite, with a fixed absolute value, and let $X/K$ be smooth,
projective and geometrically integral of dimension $n$. For a rational
codimension-$r$ cycle $Z$, a zero-curvature Green form is required to satisfy
\begin{equation}\label{eq:strong-G}
 \int_{X^{\an}\setminus |Z|^{\an}}\gamma_Z\wedge\D\eta
       =-\int_{Z^{\an}}\eta,
 \qquad \eta\in B_c^{n-r,n-r}(X^{\an}),
\end{equation}
where $B^{\bullet,\bullet}$ is Mihatsch's sheaf of ordinary actual
$\delta$-forms and
$\gamma_Z\in B^{r-1,r-1}(X^{\an}\setminus |Z|^{\an})$ has the
integrability required in this identity. The support $|Z|$ is the support of
the specified cycle after equal components have been collected. All forms,
tests and integrals in \eqref{eq:strong-G} use the original valued field $K$.

For complete intersections cut out by global functions, Mihatsch constructs
such forms explicitly \cite[Theorem 5.5]{Mihatsch}. The general projective
problem requires a global compatibility condition. The companion obstruction
theory now constructs a Chow-invariant actual Bott--Chern class and a curved
pair for every prescribed cycle \cite[Theorem 1.1]{ChowGreenPairs}. Vanishing of this obstruction
allows one to correct a curved Green form for the original cycle by a global
ordinary form. The algebraic question is therefore to identify Chow sources
on which this obstruction vanishes.

The present paper develops this question through Chow decompositions,
semistable models, and explicit geometric constructions. Its analytic
foundation is stated in \secref{sec:analytic}: the cycle obstruction is
established, while the numerical-model, Chow-covariance and height
premises are retained at their points of use. We first describe the
abelian transfer, which also provides a classical formula for the height.
The required algebraic datum is literal.
For smooth projective varieties $V,W$, a correspondence of degree $a$ from
$V$ to $W$ means an element of
$\CH^{\dim V+a}(V\times W)_\Q$; it sends codimension $r$ to codimension
$r+a$. Write $\CH^r(V)_{\Q,\homol}$ for the kernel of the geometric
cycle class map, with homological triviality imposed on every geometric
component when $V$ is disconnected.

\begin{definition}\label{def:abelian-retraction}
A smooth projective geometrically integral variety $X/k$ has a
\emph{potential abelian Chow retraction} if there are a finite extension
$L/k$, abelian varieties $A_j/L$, integers $a_j$, and rational
correspondences
\[
 u_j\in\CH^{\dim X+a_j}(X_L\times_L A_j)_\Q,
 \qquad
 v_j\in\CH^{\dim A_j-a_j}(A_j\times_L X_L)_\Q
\]
such that
\begin{equation}\label{eq:diagonal-intro}
       \Delta_{X_L}=\sum_j v_j\circ u_j
       \quad\text{in }\CH^{\dim X}(X_L\times_L X_L)_\Q.
\end{equation}
The zero-dimensional abelian variety $\Spec L$ is allowed.
\end{definition}

\paragraph{Analytic constructions used below.}
Throughout this paper, assertions about actual Green forms use the
cycle obstruction of \propref{prop:obstruction} and the
numerical-model and Chow-covariance propositions in
\secref{sec:analytic}. Assertions about the analytic height additionally
use the reciprocity and comparison theorem in \subref{sub:height-input}.
Each further model theorem is used with the hypotheses stated at its point
of application. This convention separates the geometric deductions proved
here from their companion analytic foundation.

\begin{theorem}[Chow-motivic transfer]\label{thm:main}
Use the local constructions of \secref{sec:analytic}.
Let $K/\Q_p$ be finite with its fixed absolute value. Let $X/K$ be smooth,
projective and geometrically integral, and suppose that $X$ has a potential
abelian Chow retraction. Then
\[
 \ob^r_{X,K}(z)=0
 \quad\text{for every }z\in\CH^r(X)_{\Q,\homol}
 \quad\text{and every }r.
\]
For $1\le r\le\dim X$, every specified rational cycle $Z$ representing
such a class has an ordinary actual Green form on
$X^{\an}\setminus |Z|^{\an}$ satisfying \eqref{eq:strong-G}.
For a fixed $r$, it is enough that the correspondences in
\eqref{eq:diagonal-intro} act as the identity on
$\CH^r(X_L)_{\Q,\homol}$.
\end{theorem}

\proofloc{proof:main}

We apply this theorem to three classes of varieties. The block conditions in
the following statement are geometric conditions over an algebraic closure.

\begin{theorem}[Explicit families]\label{thm:families}
Over a field $k$ of characteristic zero, each of the following varieties
has a potential abelian Chow retraction, as does every finite product of
them and every smooth projective geometrically integral variety becoming
isomorphic to such a product over a finite extension.
\begin{enumerate}[label=\textup{(\roman*)}]
\item A positive-dimensional hypersurface
\begin{equation}\label{eq:binary-family}
 \sum_{i=1}^{b}f_i(x_{i0},x_{i1})+\sum_{j=1}^{a}c_jt_j^d=0
 \quad\text{in }\PP^{2b+a-1}_k,
\end{equation}
where $d\ge2$, $2b+a\ge3$, the variable blocks are disjoint, each
$c_j\ne0$, and each binary form $f_i$ has a reduced zero divisor on $\PP^1$.
\item A cubic hypersurface in $N=5b+a\ge3$ variables, where
$0\le a\le4$, given by
\begin{equation}\label{eq:cubic-family}
       \sum_{i=1}^{b}f_i(x_{i1},\ldots,x_{i5})+g(y_1,\ldots,y_a)=0.
\end{equation}
Each $f_i=0$ is a smooth cubic threefold. For $a=0$ the last term is
omitted; for $a=1$ it is $cy_1^3$ with $c\ne0$; for $a\ge2$ it defines a smooth
cubic hypersurface in $\PP^{a-1}$.
\item The cyclic cover of $P=\PP^n_k$, $n\ge1$, defined by
\begin{equation}\label{eq:cyclic-family}
 X(F;q,m)=\Spec_P\left(\bigoplus_{j=0}^{q-1}\mathcal O_P(-jm)\right),
 \qquad u^q=F,
\end{equation}
where $q\ge2$, $m\ge1$, and $F$ is a degree-$qm$ polynomial in
$n+1=2b+a\ge2$ variables, with disjoint reduced binary blocks and
nonzero unary terms as in \eqref{eq:binary-family}.
\end{enumerate}
Over every finite $K/\Q_p$ these varieties satisfy the
conclusion of \thmref{thm:main} for every residue prime $p$.
\end{theorem}

The proofs are in \subref{proof:binary}, \subref{proof:cubic} and
\subref{proof:cyclic}, respectively. Products are handled by a tensor
identity of Chow correspondences. The cubic motive in (ii) is the
abelian-type motive of Laterveer \cite[Theorem 5]{Laterveer}; the work
needed here is its finite algebraic descent and its use in the original
valued-field theory. The construction in (i) follows the controlled join
of Shioda--Katsura \cite[Remark 1.10]{SK}.

\begin{theorem}[Normalized height formula]\label{thm:height-main}
Use the local and global constructions of \secref{sec:analytic}.
Let $F$ be a number field and let $X/F$ be smooth, projective and
geometrically integral of dimension $n\ge1$, with a potential abelian
Chow retraction. The canonical normalized analytic height of the companion
height theory is defined on
\begin{equation}\label{eq:height-domain}
 \CH^r(X)_{\Q,\homol}\times\CH^{n+1-r}(X)_{\Q,\homol},
 \qquad 1\le r\le n.
\end{equation}
For any factorization \eqref{eq:diagonal-intro} over a finite $E/F$,
\begin{equation}\label{eq:height-main}
 h_X^{\an}(z,w)=\frac1{[E:\Q]}\sum_j
       \bigl\langle u_{j*}z_E,{}^tv_{j*}w_E\bigr\rangle_{A_j,\mathrm{Kun}}.
\end{equation}
Here the classical abelian pairing has the normalization of K\"unnemann
\cite{Kun}. The two image codimensions are
$r+a_j$ and $\dim A_j+1-r-a_j$. Terms in zero or impossible homological
degrees are omitted. The value is independent of the factorization and
the finite field of definition.
\end{theorem}

\proofloc{proof:height}


% Source: papers/pairing-sections/overview.tex
\subsection{The geometric domains}

For a smooth projective variety $V/K$, write $\mathsf V(V)$ for
\[
 \CH^r(V)_{\Q,\geomhom}\subseteq\ker\ob^r_{V,K}
 \quad\text{for every }r.
\]
Here and below, the local assertions use the companion results stated in
\secref{sec:analytic}. A number-field variety has a \emph{full height
domain} if both complete groups in \eqref{eq:height-domain} are in the
domain of the companion normalized pairing for every $r$.

\begin{theorem}[Geometric domains]\label{thm:geometric-domains}
Use the analytic constructions of \secref{sec:analytic} and assume the model
inputs specified in the cited sections.
Let $K/\Q_p$ be finite. Each of the following smooth projective
geometrically integral varieties has property $\mathsf V$.
\begin{enumerate}[label=\textup{(\roman*)}]
\item A variety of dimension at most three; a geometrically uniruled
fourfold; or a variety of dimension at most five whose universal
geometric rational $\CH_0$ is supported in dimension at most three.
\item A product
\[
 T\times\prod_a C_a\times\prod_b S_b,
\]
where the $C_a$ are arbitrary smooth projective curves, the $S_b$ are
arbitrary smooth projective surfaces, and $T$ has potentially good
projective reduction. The coefficient $T$ may be a point.
\item A variety whose rational Chow motive, after a finite extension,
is a retract of a finite direct sum of Tate twists of the varieties in
\textup{(ii)}. In particular this includes products of point Hilbert
schemes of arbitrary surfaces with curves and the allowed coefficient
$T$, as well as the surface-moduli motives specified in
\secref{sec:products}.
\item A variety in one of the full-degree Chow-niveau, complete-intersection,
or flat-fibration cases of \secref{sec:niveau}, or satisfying an
all-degree numerical model criterion of \secref{sec:models}.
\end{enumerate}
For every specified geometrically homologically trivial cycle, property
$\mathsf V$ supplies a form satisfying \eqref{eq:strong-G} on its exact
support complement. For a number-field variety to which one of these
criteria applies at every finite completion, the normalized pairing has
a full height domain.
\end{theorem}

The proofs of (i), (ii)--(iii), and (iv) occupy
\secref{sec:endpoints}, \secref{sec:products}, and
\secref{sec:niveau}--\secref{sec:models}, respectively. The global
assertion follows from \lemref{lem:two-degrees}. Geometric constructions
can be base changed to all completions; a reduction condition on $T$ is
imposed separately at each finite place.

The distinction between a local degree range and a global height range
is useful even when all degrees are not covered. If $S\subseteq
\{1,\ldots,n\}$ is the set of local degrees established at every finite
place, then the full height pairs occur in
\begin{equation}\label{eq:range-intersection}
 S\cap(n+1-S).
\end{equation}
For example, small diagonal support gives local degrees $r\leq2$ or
$r\geq n-2$; for $n\geq6$ the corresponding full height degrees are
$r\leq2$ or $r\geq n-1$.

On an arbitrary smooth projective variety, the same formalism applies
to subgroups. In particular algebraically trivial cycles lie in the
local obstruction kernel in every degree. This gives the pairing on
the two algebraically trivial groups described in
\secref{sec:cycle-domains}.



\subsection{Proof strategy and relation to earlier work}

The local transfer has three steps. First, K\"unnemann's exact complex on
a strictly semistable abelian model produces a cycle extension whose
restriction to every special-fiber component is homologically trivial.
The companion numerical criterion then annihilates the actual
obstruction. Second, the diagonal identity expresses an original Chow
class as a sum of images of homologically trivial abelian classes.
Chow-kernel covariance transfers their vanishing and the full finite
degree descends it to $K$. Third, the original curved Green pair is
corrected on $X$ itself. This last step retains the original support.

For the separated-variable families, a single join step has source
$\Bl_C(A\times B)$ with a smooth codimension-two center $C$. Both the
ambient product and the center have lower-complexity Chow factorizations.
The exact blowup identity and the degree retraction therefore give a
recursive factorization for the target. For cyclic covers, the additional
map from an ordinary hypersurface is finite flat everywhere, of rank $m$.
Its full rank is the factor in the Chow identity and in the height formula.

The geometry of abelian-type motives and the classical abelian height are
inputs to this argument. K\"unnemann's height already has correspondence
functoriality and finite-extension compatibility \cite[\S1]{Kun}. Our
height formula records its realization by the companion ordinary-form
theory and the precise scalar in the explicit Chow constructions.
The companion references in \secref{sec:analytic} separate this analytic
input from the algebraic arguments proved here.

After the analytic input and abelian transfer, we treat the endpoint
degrees, small diagonal support, products, Chow niveau, and flat
fibrations. The explicit Chow identities then give separated hypersurfaces
and cyclic covers. The final sections discuss model criteria, global
families, and the normalized height formulas.

\section{The analytic constructions}\label{sec:analytic}

\subsection{Ordinary forms and the cycle obstruction}\label{sub:obstruction}

Fix $K$ as above and set $D=d'd''$. For a smooth projective $K$-variety,
the actual Bott--Chern group used in the companion theory is
\begin{equation}\label{eq:BC}
 H_{\BC}^{r,r}(X^{\an})=
 \frac{\{\omega\in B^{r,r}(X^{\an}):d'\omega=d''\omega=0\}}
 {D B^{r-1,r-1}(X^{\an})}.
\end{equation}
The real vector-space structure of this group is retained, and the
cycle map below is $\Q$-linear.

\begin{proposition}[Actual cycle obstruction]\label{prop:obstruction}
There is a Chow-invariant map
\[
 \ob^r_{X,K}:\CH^r(X)_\Q\longrightarrow H_{\BC}^{r,r}(X^{\an}).
\]
For a specified cycle $Z$, the construction gives an ordinary actual
form $G_Z$ on its exact support complement and a closed global ordinary
form $\Omega_Z$ such that
\begin{equation}\label{eq:curved}
       D[G_Z]+\delta_Z=[\Omega_Z],
       \qquad \ob^r_{X,K}([Z])=[\Omega_Z].
\end{equation}
The equality is tested against every complementary compact actual form,
with the integrability supplied by the construction. If
$\ob^r_{X,K}([Z])=0$, there is a global
$\Psi\in B^{r-1,r-1}(X^{\an})$ with $D\Psi=\Omega_Z$, and
$G_Z-\Psi$ satisfies \eqref{eq:strong-G} on the same complement.
\end{proposition}

\begin{proof}
For integral $X$, the Chow map and curved pair are
\cite[Theorem 1.1]{ChowGreenPairs}; the theorem applies to every
prescribed cycle over the original valued field, in every positive
codimension up to the dimension of $X$. A smooth projective variety
is the disjoint union of its finitely many integral connected
components, so the construction applies componentwise as well.
The definition of the quotient in \eqref{eq:BC} supplies the global
primitive when the curvature class is zero. Subtracting this
primitive and applying actual Stokes gives \eqref{eq:strong-G} on
the same complement.
\end{proof}

The current notation in \eqref{eq:curved} expresses the all-actual-test
integral identity of the supplied ordinary forms. The finite numerical
criterion and correspondence covariance below supply additional
vanishing and functoriality statements, with their own proofs.

\subsection{Numerical witnesses and Chow covariance}\label{sub:covariance}

We state the numerical input in the finite constant-field case used in
this paper. A pure operational representative of a codimension-$r$
cycle on a regular model means a rational virtual vector bundle whose
Chern character has that cycle in degree $r$ and vanishes in all other
degrees. Such a representative is part of the companion construction.

\begin{proposition}[Finite numerical-model criterion]\label{prop:numerical}
Let $V/K$ be smooth projective and integral, let $1\le r\le\dim V$,
and let $z\in\CH^r(V)_\Q$. Suppose that a smooth integral finite
constant-field cover $V'\to V$, of complete native degree $e>0$, has a
regular projective model carrying a pure operational representative of
the pulled-back class $z$. Suppose that its degree-$r$ character pairs
to zero with every integral vertical $r$-cycle, with all
scheme-theoretic fiber multiplicities and residue degrees retained.
Then $\ob^r_{V,K}(z)=0$.
\end{proposition}
\begin{proof}
Let $h:V'\to V$ be the supplied cover and let $\xi$ be its numerical
model representative. The numerical model-Chern detector makes
$\operatorname{ch}_r(\bar\xi)=0$ as an actual form over $K$.
Choose a model curvature $\Omega$ representing $\ob^r_{V,K}(z)$.
The two characters $h^*\Omega$ and $\operatorname{ch}_r(\bar\xi)$
have the same generic Chow class; the generic-Chow-zero construction
therefore supplies a global actual $\beta$ with $D\beta=h^*\Omega$.

Both generic varieties are smooth, hence normal. Full finite actual
trace, \cite[\texttt{r16ft:full-trace}]{InternalCovariance}, gives
\[
 \Psi=e^{-1}\Tr_h\beta,\qquad
 D\Psi=e^{-1}\Tr_h(h^*\Omega)=\Omega.
\]
Thus $\ob^r_{V,K}(z)=0$. This is also the conclusion of
\cite[\texttt{r12a:finite-numerical-cover}]{InternalNumerical},
whose earlier proof traces the original scalar--curvature generators
along composite finite maps. If a model is over a larger constant
field, viewing it over $K^\circ$ multiplies vertical numerical degrees
by the residue degree and retains ramification multiplicities.
Zero remains zero, and the trace denominator is the full native $e$.
\end{proof}

This is the finite-cover case of \cite{InternalNumerical}, Theorem
\texttt{r12a:finite-numerical-cover}. It applies as well when the
constant field of $V$ is a finite extension of $K$: the model and its
numerical degrees are then regarded over $K^\circ$.

Write $\mathcal K_K^r(V)=\ker\ob^r_{V,K}$. The following covariance
is an assertion about Chow classes.

\begin{proposition}[Native Chow-kernel covariance]\label{prop:covariance}
In the fixed $K$-theory, rational Chow correspondences between smooth
projective varieties preserve $\mathcal K_K^*$, with their usual
codimension shifts. In particular proper Chow pushforward preserves
the kernel. For a finite extension $L/K$, the projection
$\pi:V_L\to V$ satisfies
\[
          \pi_*\pi^*z=[L:K]z.
\]
Thus native-$K$ obstruction zero of $\pi^*z$ implies obstruction zero
of $z$. Finite actual trace and integration use the full degree
$[L:K]$ and the absolute value extending that of $K$.
\end{proposition}
\begin{proof}
The actual obstruction is a graded-ring map and commutes with
algebraic pullback. For the projection $\pi:T\times\PP^N\to T$,
write $a=\sum_{j=0}^N h^j\pi^*a_j$ by the Chow projective-bundle
formula. The all-actual projective-bundle operator commutes with $D$
and satisfies the target-factor formula and
$I_\pi(h^j)=0$ for $j<N$, $I_\pi(h^N)=1$.
These are \cite[\texttt{r16out:line}, \texttt{r16out:bundle}]{InternalProjective}.
Consequently $I_\pi\ob(a)=\ob(a_N)=\ob(\pi_*a)$, so this
projective push preserves the zero obstruction class.

For a smooth closed immersion, the deformation-to-the-normal-cone
proof in \cite[\texttt{r16gi:detection}, \texttt{r16gi:kernel}]{InternalProjective}
uses the full Bott--Chern projective-bundle formula and smooth-center
blowup injectivity. Its special restrictions are zero and a Chern
polynomial times the pulled-back source class. A zero source
obstruction therefore gives zero on the general fibre, which is
the obstruction of the algebraically pushed class. A proper morphism
between smooth projective varieties factors as its graph immersion
into a product with projective space followed by that product's
projection. Combining the two statements proves stability for proper
Chow pushforward; this is \cite[\texttt{r16cov:proper-push}]{InternalCovariance}.

For a correspondence $\Gamma$, apply pullback and the ring property
to $\Gamma\cdot p_1^*z$, then proper kernel stability to
$\Gamma_*z=p_{2*}(\Gamma\cdot p_1^*z)$. This proves the assertion
with the usual codimension shifts. The argument applies on each
integral component of a smooth projective product. Finally,
$\pi_*\pi^*z=[L:K]z$ is the algebraic finite-flat degree identity.
Kernel stability and $\Q$-linearity give descent. The finite trace
and every integration identity retain this full native degree.
\end{proof}

The source is \cite{InternalCovariance}, Theorem
\texttt{r16cov:correspondences}, together with the full finite trace
theorem FT16. If $L/K$ has ramification index $e(L/K)$ and residue
degree $f(L/K)$, then
\begin{equation}\label{eq:local-normalization}
 [L:K]=e(L/K)f(L/K),\qquad
 c_L=-\log|\varpi_L|=c_K/e(L/K).
\end{equation}

There is a useful distinction in this convention. If $V$ is defined
over $L$, its native obstruction $\ob_{V,K}$ is calculated on $V$
regarded as a $K$-scheme. A correspondence over $L$ is inserted into
the open-and-closed diagonal factor
\[
       V\times_L W\hookrightarrow V\times_K W.
\]
Composition on this factor is composition over $L$. This realizes an
$L$-Chow identity inside the original $K$-theory without changing its
coefficients.

\subsection{Heights on the admitted Chow domain}\label{sub:height-input}

For number fields, use the standard absolute values: the finite ones
extend $|p|=p^{-1}$ and the archimedean ones are the usual absolute
values, with the companion product-formula weights. For a
smooth projective geometrically integral variety $X/F$ of dimension $n$,
let $\mathcal H^r(X)$ be the subgroup
of $\CH^r(X)_{\Q,\homol}$ whose localizations at all finite places
have zero actual obstruction.

\begin{theorem}[Admitted-domain reciprocity and comparison]\label{hyp:height-input}
For $1\le r\le n$, there is a canonical normalized rational-bilinear real pairing
\[
 h_X^{\an}:\mathcal H^r(X)\times\mathcal H^{n+1-r}(X)\longrightarrow\R.
\]
It includes the archimedean terms, local residue identities, global
reciprocity, symmetry in complementary codimensions, and descent to rational equivalence. On admitted classes
it is invariant under finite number-field extension. For an actual
Chow correspondence $v:A\rightsquigarrow X$ between smooth projective
varieties, with classes in the indicated admitted domains and
complementary codimensions, its adjunction formula is
\begin{equation}\label{eq:height-adjunction}
       h_X^{\an}(v_*a,w)=h_A^{\an}(a,{}^tv_*w).
\end{equation}
For an abelian variety $A/E$, the comparison theorem gives
\begin{equation}\label{eq:abelian-comparison}
       h_A^{\an}(a,b)=[E:\Q]^{-1}\langle a,b\rangle_{A,\mathrm{Kun}}.
\end{equation}
\end{theorem}
\begin{proof}
For a class in $\mathcal H$, \propref{prop:obstruction} supplies
exact-support strong Green forms at every finite completion.
The local construction in \cite[\texttt{r16h:curved-slice},
\texttt{r16h:global-height}]{InternalHeight} gives three facts.
Disjoint-polar Stokes proves reciprocity and independence of the
Green choices. Spreading disjoint cycle representatives to good
models gives zero local value outside finitely many places. Finally,
the supported point-slice identity on the normal carriers of a
properly moved higher-Chow one-chain gives the principal variation
as minus the logarithm of one algebraic norm in the ground number
field. Its refined point coefficients include all lengths and
residue-field degrees.

Sum with the product-formula weights. The sum is finite; the product
formula kills the principal variation and hence gives descent in
each argument to rational equivalence. Local reciprocity proves
symmetry. Moving finitely many cycles into a common disjoint position
proves bilinearity. The native local finite-degree formula and
$\sum_{w\mid v}[E_w:F_v]=[E:F]$ give invariance under number-field
extension with the absolute normalization. These are precisely the
construction and independence claims of H16.

By \propref{prop:covariance}, correspondence actions preserve these
domains. Resolve a prime correspondence to a smooth projective
variety with maps to its two factors. The simultaneous moving lemma
\cite[\texttt{r16ha:moving}]{InternalHeight} makes the two refined
pullbacks proper and disjoint. Apply the proper height projection
formula to each map, with local symmetry between the two
applications. Their algebraic refined pushforwards are the actions
of $v$ and ${}^tv$, which gives \eqref{eq:height-adjunction}; this
is \cite[\texttt{r16ha:adjunction}]{InternalHeight}.

For an abelian variety, after finite extension K\"unnemann's
model construction supplies numerically admissible lifts
\cite[Theorem 1.6]{Kun}.
The framed local comparison identifies each Green value with its
finite Euler intersection, multiplied by the native valuation
constant. The weighted sum, including the matching archimedean
terms, is the arithmetic model height divided by the absolute field
degree; see \cite[\texttt{r16mm:local}, \texttt{r16mm:global}]{InternalHeight}.
K\"unnemann's unnormalized height descends across $E'/E$ by
$[E':E]^{-1}$ \cite[\S1.4, equation (7), and Lemma 1.5]{Kun},
while the normalized analytic pairing is invariant.
Since $[E':\Q]=[E':E][E:\Q]$, the comparison descends to
\eqref{eq:abelian-comparison}, as proved in
\cite[\texttt{r16mm:abelian}]{InternalHeight}.
\end{proof}

The three assertions use H16, HA16 and MM16 with their stated domains.
For a general variety the model comparison requires supplied regular
models and numerically admissible lifts; it does not assert their
existence for every cycle. The domain $\mathcal H$ is the intersection
of all local obstruction kernels. Replacing it by the kernel of a
chosen finite collection of numerical tests requires an additional
comparison, and enlarging it to every homologically trivial class
requires the corresponding local obstruction vanishing.

\section{Abelian varieties and Chow-motivic transfer}\label{sec:transfer}

\subsection{The local abelian model}

Let $R$ be a complete discrete valuation ring of characteristic-zero
fraction field $F$ and finite residue field $k$. Let $\mathcal P/R$ be
regular, projective and strictly semistable of relative dimension $d$,
with special fiber $Y=\bigcup_iD_i$. For a smooth projective
$k$-variety $D$, put
\[
 A^q(D)=\operatorname{im}\bigl(\CH^q(D)_\Q\to
 H^{2q}_{\textnormal{ét}}(D_{\bar k},\Q_\ell(q))\bigr),
 \qquad \ell\ne\operatorname{char}k.
\]
Using the signed restriction $\rho$ and Gysin $\gamma$ maps on the
component intersections, define
\begin{align*}
 A^q(Y)&=\ker\bigl(A^q(Y^{(1)})\xrightarrow{\rho}A^q(Y^{(2)})\bigr),\\
 A_s(Y)&=\operatorname{coker}\bigl(
 A^{d-s-1}(Y^{(2)})\xrightarrow{\gamma}A^{d-s}(Y^{(1)})\bigr).
\end{align*}
Here $Y^{(a)}$ is the disjoint union of the $a$-fold intersections.

For an abelian variety over $F$, K\"unnemann constructs such a model
after a finite extension and proves that the natural homology map
$A_s(Y)\to H^{\textnormal{ét}}_{2s}(Y_{\bar k},\Q_\ell(-s))$
is injective and that
\begin{equation}\label{eq:kun-complex}
 A_{d+1-r}(Y)\xrightarrow{i^*i_*}A^r(Y)
                  \xrightarrow{\cap[Y]}A_{d-r}(Y)
\end{equation}
is exact at the middle term \cite[Theorems 5.6 and 6.6]{Kun}.
The complete-DVR clause of the correction in \cite[\S5.8]{Kun}
is the one used here.

\begin{lemma}\label{lem:model-correction}
Assume the injectivity and exactness just stated. Every
$z\in\CH^r(\mathcal P_F)_{\Q,\homol}$ has an extension
$\Theta_0\in\CH^r(\mathcal P)_\Q$ whose refined restriction to
each $D_i$ is geometrically homologically trivial.
\end{lemma}

\begin{proof}
Choose a cycle extension $\Theta$ and let $a\in A^r(Y)$ be its
component restrictions. The class $a\cap[Y]$ is the cycle-image class
of the specialization of $z$. Its geometric homology class is zero
because specialization commutes with the cycle class map. Injectivity
gives $a\cap[Y]=0$ in $A_{d-r}(Y)$, and exactness of
\eqref{eq:kun-complex} supplies $b\in A_{d+1-r}(Y)$ with
$i^*i_*b=a$.

The map $\CH_{d+1-r}(Y)_\Q\to A_{d+1-r}(Y)$ is surjective.
Lift $b$ to $\beta$ and set $\Theta_0=\Theta-i_*\beta$. Its generic
restriction is $z$, and its tuple of component cycle classes is zero.
Since $A^r(Y)$ is a subspace of the direct sum of the component
cycle-image groups, each component restriction has zero class.
\end{proof}

\begin{proposition}\label{prop:abelian-zero}
For every abelian variety $A/K$ and every $r$,
$\CH^r(A)_{\Q,\homol}\subset\mathcal K_K^r(A)$. The same
holds for an abelian variety over a finite $L/K$ regarded as a
$K$-scheme, with homological triviality on every geometric component.
\end{proposition}

\begin{proof}
Choose the finite extension and the abelian model supplied above.
By \lemref{lem:model-correction}, the pulled-back cycle has an
extension with homologically trivial component restrictions. These
restrictions are numerically trivial. Its pure operational
representative therefore has zero pairing with every integral
vertical $r$-cycle: such a cycle lies in a special-fiber component,
and the pairing is the numerical pairing of the refined restriction
there. \propref{prop:numerical} gives the desired native obstruction
zero.

If $A$ is over $L$, choose its semistable extension over a finite
$M/L$ and regard the model over $K^\circ$. Residue degrees multiply
the numerical pairings, and ramification is included in the
scheme-theoretic multiplicities. Their zero values are unchanged.
The same criterion applies to the finite $K$-morphism $A_M\to A$.
The degree-zero and out-of-range homological Chow groups vanish.
\end{proof}

\subsection{Proof of the transfer theorem}\label{proof:main}

Let $L/K$ and $(u_j,v_j)$ be a retraction as in
\eqref{eq:diagonal-intro}, and let $z\in\CH^r(X)_{\Q,\homol}$.
Every $u_{j*}z_L$ is geometrically homologically trivial.
By \propref{prop:abelian-zero}, it lies in the native group
$\mathcal K_K^{r+a_j}(A_j)$. Applying \propref{prop:covariance} to
$v_j$ and summing the diagonal identity gives
\[
               z_L=\sum_jv_{j*}u_{j*}z_L\in\mathcal K_K^r(X_L).
\]
As explained in \subref{sub:covariance}, this is an identity in the
original $K$-theory. The finite projection $\pi:X_L\to X$ gives
$\pi_*z_L=[L:K]z$. Kernel covariance and rational division by the
complete degree imply $\ob^r_{X,K}(z)=0$.

For the specified cycle $Z$, choose the pair $(G_Z,\Omega_Z)$ of
\propref{prop:obstruction}. Its vanishing class supplies a global
$\Psi$ with $D\Psi=\Omega_Z$. Thus $\gamma_Z=G_Z-\Psi$ has
the required support, integrability and test identity. The same
argument proves the fixed-degree version, since only the action
of the diagonal on $z_L$ was used. This proves \thmref{thm:main}.


% Source: papers/pairing-sections/endpoints_support.tex
\section{Endpoint degrees and small diagonal support}\label{sec:endpoints}

\subsection{Chow transfer and complementary degrees}

\begin{lemma}[Transfer through a Chow diagonal]\label{lem:general-transfer}
Suppose that over a finite extension $L/K$ there are smooth projective
varieties $V_j$ and correspondences of opposite degrees satisfying
\[
 \Delta_{X_L}=\sum_jv_j\circ u_j
 \quad\text{in }\CH^{\dim X}(X_L\times_LX_L)_\Q.
\]
Fix $r$. If every homologically trivial image $u_{j*}z_L$ is in the
native obstruction kernel of $V_j$, then
$z\in\cK_K^r(X)$. In particular $\mathsf V(V_j)$ for every $j$
implies $\mathsf V(X)$.
\end{lemma}

\begin{proof}
Cycle classes commute with correspondences, so the image classes are
homologically trivial. Kernel covariance sends $v_{j*}u_{j*}z_L$ into
$\cK_K^r(X_L)$. Sum the diagonal and push down by $X_L\to X$.
The result is $[L:K]z$, and rational division gives the conclusion.
All correspondences over $L$ are interpreted on the diagonal
open-and-closed factor of the product over $K$, as in
\subref{sub:covariance}. Finally apply the curved-pair theorem to the
specified representative of $z$.
\end{proof}

\begin{lemma}[The two-degree criterion]\label{lem:two-degrees}
Let $X/F$ have dimension $n$, and let $S\subseteq\{1,\ldots,n\}$.
Suppose that at every finite place $v$, all geometric homological
classes in degrees $r\in S$ have zero native obstruction. Then
$\cH^r(X)=\CH^r(X)_{\Q,\geomhom}$ for $r\in S$. The full height
pairing in degrees $(r,n+1-r)$ is therefore defined whenever
$r\in S\cap(n+1-S)$.
\end{lemma}

\begin{proof}
Geometric cycle classes commute with passage to $F_v$. Thus every
global homologically trivial class in an asserted degree belongs to
each local kernel defining $\cH^r(X)$. The opposite inclusion is in
the definition of $\cH^r(X)$. Apply the companion pairing to both
resulting equalities.
\end{proof}

\begin{proposition}[Endpoint input]\label{prop:endpoint-input}
For a smooth projective $n$-fold over the original $K$, geometric
homological divisors, zero-cycles and one-cycles have zero obstruction.
Thus $\mathsf V(X)$ holds if $n\leq3$.
\end{proposition}

The divisor and zero-cycle assertions are the endpoint results in
\cite{InternalEndpoints}. For one-cycles, the precise input is the
regular-model curve correction followed by dominant obstruction
descent, also recorded there. Its proof proceeds as follows. Choose a
projective strictly semistable alteration $V\to X$, over a finite
constant extension, of the same dimension as $X$. The pullback of a
homologically trivial one-cycle remains a one-cycle. The curve
correction produces a rational virtual bundle on the regular model
whose generic degree-$(n-1)$ character is that class and whose degree
on every vertical $(n-1)$-cycle is zero. The numerical-model criterion
annihilates its obstruction; dominant descent gives the assertion on
$X$. This uses the full degree of the alteration and leaves the final
curved-pair correction on $X$.

\subsection{Two orientations of the diagonal}\label{sub:small-support}

\begin{definition}\label{def:small-support}
A rational small-support diagonal for $X$, of dimension $n$, is an
identity
\begin{equation}\label{eq:small-diagonal}
 \Delta_X=A+B,\qquad
 \operatorname{Supp}A\subseteq D\times X,\quad
 \operatorname{Supp}B\subseteq X\times T,
\end{equation}
in rational Chow, where $D$ is a divisor and $\dim T\leq3$.
The identity may be supplied after a finite extension of the ground
field. The supports belong to actual cycle representatives of $A,B$.
\end{definition}

\begin{theorem}[Small diagonal support]\label{thm:small-support}
If $X$ has a small-support diagonal, its complete geometric homological
Chow groups are in the local obstruction kernel for
\begin{equation}\label{eq:small-support-range}
 r\leq2\quad\text{or}\quad r\geq n-2.
\end{equation}
It has property $\mathsf V$ if $n\leq5$. Over a number field, the
same geometric hypothesis gives a full height domain if $n\leq5$;
for $n\geq6$ it gives the full pairs with $r\leq2$ or $r\geq n-1$.
\end{theorem}

\begin{proof}
Resolve the components of $D$ and $T$, with maps $a_i:D_i\to X$ and
$b_j:T_j\to X$. A cycle supported on the image of a proper surjection
can be lifted rationally: choose a closed point of the fiber above its
generic point, take its closure, and divide by the full generic
degree. This also applies to cycles contained in exceptional images.
It lifts \eqref{eq:small-diagonal} to correspondences
\[
 U_i\in\CH^{n-1}(D_i\times X)_\Q,\qquad
 V_j\in\CH^{t_j}(X\times T_j)_\Q,\quad t_j=\dim T_j\leq3.
\]
The identity and its transpose act as
\begin{align}
 z&=\sum_iU_{i*}a_i^*z+\sum_jb_{j*}V_{j*}z,
       \label{eq:small-forward}\\
 z&=\sum_i a_{i*}{}^tU_{i*}z+\sum_j{}^tV_{j*}b_j^*z.
       \label{eq:small-transpose}
\end{align}
All intermediate classes are geometrically homologically trivial.
In \eqref{eq:small-forward}, the class on $D_i$ has dimension
$n-1-r\leq1$ when $r\geq n-2$. In
\eqref{eq:small-transpose}, it has codimension $r-1\leq1$ when
$r\leq2$. These are the endpoint degrees. Every feasible class on
$T_j$ is an endpoint class because $t_j\leq3$; impossible degrees
and homologically trivial rank classes are zero. Apply kernel
covariance to the indicated orientation. If the identity is defined
after a finite extension, use \lemref{lem:general-transfer} to descend.
The two local ranges cover every degree for $n\leq5$. Intersecting
them with their complementary ranges proves the height assertions
by \lemref{lem:two-degrees}.
\end{proof}

\begin{proposition}[Universal zero-cycle support]\label{prop:universal-support}
Suppose that a closed subset $T\subset X_{\bar k}$, of dimension at
most $s$, surjects onto $\CH_0(X_\Omega)_\Q$ for every algebraically
closed extension $\Omega/\bar k$. Then after a finite extension of
$k$ there is an identity \eqref{eq:small-diagonal} with $\dim T\leq s$.
\end{proposition}

\begin{proof}
Apply the support hypothesis to the generic point over
$\overline{\bar k(X)}$. A rational-equivalence certificate uses only
finitely many points, curves and functions, so it is defined over a
finite extension of $\bar k(X)$. Norming down and dividing by its
degree gives the supported generic-point relation over $\bar k(X)$.
Close the cycle in $X_{\bar k}\times T$ and spread the relation over
an open subset of the first factor. Chow localization puts the
remaining difference from $\Delta_X$ over its closed complement;
contain this complement in a divisor. Finally descend the finite
cycle-and-function certificate and the supports to a finite
extension of $k$.
\end{proof}

\begin{corollary}[Rationally connected fibrations]\label{cor:rc-fibrations}
Let $f:X\to B$ be a surjective morphism of smooth projective varieties
in characteristic zero. Suppose that over a nonempty open of $B$ it
is smooth with rationally connected geometric fibers. If $\dim B\leq3$,
then \thmref{thm:small-support} applies to $X$. In particular it gives
all degrees if $\dim X\leq5$. Geometrically rationally connected
varieties, including smooth Fano varieties, are examples with
zero-dimensional support.
\end{corollary}

\begin{proof}
The closure of a closed point of the generic fiber is a multisection
$T\subset X$ of dimension $\dim B$. It is proper and surjective over
$B$. After any algebraically closed extension, move a zero-cycle
into the smooth rationally connected locus. Each of its points is
rationally equivalent, inside its fiber, to a point of $T$ in the
same fiber. Rational connectedness gives this equality by chains of
rational curves. The multisection therefore supplies universal
geometric rational zero-cycle support. Apply
\propref{prop:universal-support} and \thmref{thm:small-support}.
\end{proof}

\subsection{Projective bundles and dominant images}

\begin{proposition}\label{prop:closures}
Property $\mathsf V$ is preserved by projective bundles, by blowups
along smooth centers when both the ambient variety and the center
have $\mathsf V$, and by smooth projective dominant morphism images.
The assertions allow a preliminary finite extension.
\end{proposition}

\begin{proof}
The Chow-motive projective-bundle and blowup decompositions are
finite direct sums of Tate twists of the indicated smooth varieties.
Apply \lemref{lem:general-transfer}. For a projective dominant
morphism $f:V\to X$ with relative dimension $e$, choose a relatively
ample class $a$ whose top degree on the generic fiber is $d>0$.
The correspondence $d^{-1}f_*(a^e\cdot-)$ is a left inverse to
$f^*$ by the projection formula. Its action gives a rational Chow
retraction and hence the same transfer. The formula can be checked
after product with any smooth projective test variety, so it is a
Chow correspondence identity. Finite descent is as before.
\end{proof}

\begin{corollary}\label{cor:uniruled}
A geometrically uniruled smooth projective fourfold has $\mathsf V$.
So does any smooth projective dominant morphism image of a
projective-bundle tower over a variety of dimension at most three,
or over one of these uniruled fourfolds.
\end{corollary}

\begin{proof}
After a finite extension, a uniruled fourfold is dominated rationally
by $B\times\PP^1$ for a smooth projective threefold $B$. Resolve this
map in characteristic zero by smooth blowups of the four-dimensional
source. Every nontrivial center has dimension at most two, hence has
$\mathsf V$. The source and its resolution have $\mathsf V$ by
\propref{prop:closures}. The resolved dominant morphism transfers it
to $X$. The remaining assertions follow successively from projective
bundle and dominant-morphism transfer.
\end{proof}



% Source: papers/pairing-sections/products_moduli.tex
\section{Products of curves and surfaces and their Chow summands}\label{sec:products}

\subsection{The product-model input}

The product theorem uses a numerical calculation on a particular
semistable model. We state its algebraic content before applying the
companion comparison with actual-form obstructions.

Let $R$ be a complete discrete valuation ring with finite residue field.
Take regular projective reduced strictly semistable models of finitely
many curves and surfaces. Fix orders on their components and form the
ordered staircase resolution of their product. Write its special fiber
as $D=\bigcup_iD_i$, with $\dim D_i=d$, retaining each connected
intersection separately. For an oriented edge $i\to j$, let
$j_i:E_{ij}\to D_i$ and $j_j:E_{ij}\to D_j$ be the inclusions. Set
\[
 (\rho x)_{ij}=j_j^*x_j-j_i^*x_i,\qquad
 \sigma y=(j_{i*}y,-j_{j*}y),\qquad A=\sigma\rho.
\]
These are cohomological actions of actual rational Chow correspondences.
The off-diagonal entries of $A$ have codimension $d+1$; its diagonal
entry on $D_i$ is
\[
 (\Delta_i)_*c_1(N_{D_i/\mathcal D})
       =-(\Delta_i)_*\sum_{j\ne i}[E_{ij}].
\]

\begin{proposition}[Companion product-model theorem]\label{prop:product-model}
For this ordered model there is an invertible prime $\ell$ and a
matrix of rational Chow correspondences $G$, with entries of
codimension $d-1$, for which
\begin{equation}\label{eq:inner-inverse}
 A^H G^H A^H=A^H
\end{equation}
on full geometric $\ell$-adic cohomology, in every degree and with
its Tate twist. This remains true after product with a smooth
projective $R$-scheme, by replacing $A,G$ with the exterior products
of their correspondences and the coefficient's diagonal.
\end{proposition}

The proof is \cite[model splitting theorem]{InternalProducts}, whose
actual deletion and residual-core constructions are included in the
proof dossier. The essential structure is as follows. Each component
is a tower of blowups of closed strata in the full reduced boundary
of a product of smooth curve and surface components. For a subset
$I$ of the surface indices, coordinate deletion gives an actual
proper birational map
\begin{equation}\label{eq:deletion-map}
 f_I:D_{\boldsymbol i}\longrightarrow
 \left(\prod_{j\in I}X_{j,i_j}\right)\times W_{-I,\boldsymbol i_{-I}},
\end{equation}
where $W_{-I}$ is the staircase component of the remaining factors.
The refined surface projectors $p_j$ onto $H^2_{\rm tr}$ annihilate
cohomology supported on the boundary curves. Consequently the transported
projectors
\[
 P_I=f_I^*((\boxtimes_{j\in I}p_j)\boxtimes1)f_{I*}
\]
satisfy $P_I=\prod_{j\in I}P_{\{j\}}$, commute with the incidence,
and do so on exceptional summands as well. The proof of commutation
uses support localization for the kernel of a birational pull-push
operator in \eqref{eq:deletion-map}.

The simultaneous sectors retain the tensor factors $t_2$ selected by
these projectors. The residual projector annihilates each generator
which still contains a full-surface $t_2$ factor; the deletion proof
establishes this equality in relative Chow before passing to
localization. The surviving blowup and localization generators
in the complementary factor are of abelian type. Thus the incidence
operator on a sector is the identity on its retained transcendental
factor tensored with an operator on one abelian carrier. Numerical
semisimplicity and the one-prime comparison for that carrier give its
inner inverse, which is tensored back with the retained factor.
Finite-field choices descend by normalized averaging: the equation
\eqref{eq:inner-inverse} is linear in $G$. This is the algebraic
product-model statement used here.

Two consequences of \eqref{eq:inner-inverse} explain its role in the
height problem. Let $\mathscr A^u$ be the space of actual algebraic
cycle classes on the disjoint union of the $D_i$, and let
$p:\mathscr A^u\twoheadrightarrow N^u(D^{(0)})_\Q\otimes\Q_\ell$
be the quotient to numerical classes. If $p(x)$ is in the numerical
kernel of $A$, then
\begin{equation}\label{eq:numerical-correction}
 x^\flat=x-G^HA^Hx
\end{equation}
has the same numerical image and belongs to $\ker A^H$. Also,
if an algebraic class $\alpha$ is in $\im A^H$, then $G^H\alpha$
is algebraic and $A^HG^H\alpha=\alpha$. Hence
\begin{equation}\label{eq:image-saturation}
 \mathscr A^{u+1}\cap\im A^H=A^H\mathscr A^u.
\end{equation}
The numerical-kernel lifting and actual-span saturation are the two
inputs to the companion model criterion \cite{InternalModel}.

\begin{theorem}[Arbitrary curve and surface products]\label{thm:products}
Let $X_1,\ldots,X_m/K$ be arbitrary smooth projective geometrically
integral curves or surfaces. Let $T/K$ be smooth projective
geometrically integral with potentially good projective reduction. Then
\[
 X=T\times\prod_{a=1}^mX_a
\]
has $\mathsf V$. The assertion holds over every finite extension of
$K$ for every newly specified geometrically homologically trivial
cycle there. Empty products and a point as coefficient are permitted.
\end{theorem}

\begin{proof}
After a common finite extension choose projective strictly semistable
alterations of the low-dimensional factors and a smooth projective
model of $T$. Semistable repair after any intermediate ramified
extension is part of the ordered-model construction in
\cite{InternalProducts}. Apply \propref{prop:product-model} to the
resulting witness. Weight--monodromy for curves and surfaces, the
purity of the good coefficient, and K\"unneth give the odd-weight
condition in the companion model criterion. Equations
\eqref{eq:numerical-correction} and \eqref{eq:image-saturation}
supply its two exactness conditions. The pulled-back obstruction
therefore vanishes. Dominant and finite-degree descent give zero
obstruction over $K$; the original curved pair gives the requested
form.

For an arbitrary finite $K'/K$, the factors still have the same
dimensions and the coefficient remains potentially good. Reapply
the construction over $K'$ to its newly specified cycle. This proves
the assertion for new classes, as well as for pullbacks of old ones.
\end{proof}

Over a number field, the curves and surfaces are unrestricted. If a
coefficient $T$ is present, it is required to have potentially good
projective reduction at every finite place. Applying
\thmref{thm:products} at all completions and using
\lemref{lem:two-degrees} gives the full height in every complementary
degree, also after every finite number-field extension.

\subsection{Abelian factors and the surface-generated Chow domain}

\begin{lemma}\label{lem:abelian-curve-source}
Over a characteristic-zero field, an abelian variety becomes a
rational Chow-motive retract of a product of curves after a finite
extension.
\end{lemma}

\begin{proof}
For an abelian variety of dimension $g\geq2$, choose a smooth
complete-intersection curve and a point $c_0$ on it after a finite
extension. Translate by $-c_0$ and call the resulting pointed curve
$C$; choose the complete intersection so that $C-C$ generates $A$.
The addition map $C^g\to A$ is dominant. Indeed,
the successive sums of $C$ increase in dimension until they fill
$A$: stabilization at a proper subvariety would make that subvariety
stable under the group generated by $C$. Since source and target
have dimension $g$, the map is generically finite, of some degree
$d>0$. Its graph satisfies $f_*f^*=d\Delta_A$, and
$d^{-1}f_*$ is a Chow retraction. Dimension one is immediate;
dimension zero gives the unit motive.
\end{proof}

\begin{corollary}\label{cor:surface-envelope}
Suppose that after a finite extension the rational Chow motive of
$X$ is a direct summand of a finite direct sum of Tate twists of
products of curves and surfaces. Then $\mathsf V(X)$ holds. It also
holds for $T\times X$ with a potentially good projective coefficient.
These conclusions include products with arbitrary abelian factors.
\end{corollary}

\begin{proof}
Choose the inclusion and retraction correspondences. Their composites
sum to the diagonal. Tensor with the coefficient's diagonal when
needed, apply \thmref{thm:products} to each intermediate variety,
and use \lemref{lem:general-transfer}. An abelian factor may first be
replaced by its curve-product retraction from
\lemref{lem:abelian-curve-source}.
\end{proof}

\subsection{Hilbert schemes of points on arbitrary surfaces}

Let $S$ be a smooth projective geometrically integral surface and
let $H=S^{[m]}$. For a
partition $\lambda=1^{a_1}\cdots m^{a_m}$, put
$l(\lambda)=\sum_i a_i$, $s_\lambda=m-l(\lambda)$, and
$G_\lambda=\prod_i\mathfrak S_{a_i}$. On the smooth power
$S^{l(\lambda)}$, use the rational projector
\[
 e_\lambda=|G_\lambda|^{-1}\sum_{g\in G_\lambda}\Gamma_g.
\]
The theorem of de Cataldo--Migliorini
\cite[Theorem 6.2.1 and Remark 6.2.2]{DCM} gives
\begin{equation}\label{eq:hilbert-motive}
 h(S^{[m]})\simeq\bigoplus_{\lambda\vdash m}
       (h(S^{l(\lambda)}),e_\lambda)\otimes\LL^{s_\lambda}.
\end{equation}
We use the covariant twist convention
$\CH_a(h(V)(s))=\CH_{a-s}(V)_\Q$.

\begin{theorem}\label{thm:hilbert}
For arbitrary smooth projective geometrically integral surfaces
$S_i$, curves $C_j$, and
positive integers $m_i$, the variety
\[
 T\times\prod_j C_j\times\prod_i S_i^{[m_i]}
\]
has $\mathsf V$ if $T$ is smooth projective geometrically integral
and has potentially good projective reduction.
Over a number field, it has a full height domain when this reduction
condition holds at every finite place.
\end{theorem}

\begin{proof}
Tensor the actual incidence correspondences in
\eqref{eq:hilbert-motive}. For a tuple of partitions they factor
through $T\times\prod_j C_j\times\prod_i S_i^{l(\lambda_i)}$.
If $s=\sum_i s_{\lambda_i}$, this intermediate variety has
dimension $\dim X-2s$ and an original codimension-$r$ class maps
to codimension $r-s$. Its homological class is zero. Apply
\thmref{thm:products} in every feasible degree and then
\lemref{lem:general-transfer}. The group in degree zero has no
homological kernel, and out-of-range groups vanish. The global
assertion follows at all completions. The correspondences commute
with base change, so the same proof applies to all new classes
after finite extension.
\end{proof}

\subsection{Finite specialization of Chow certificates}

\begin{lemma}\label{lem:specialize-certificate}
Let $X/k$ be smooth projective in characteristic zero. Suppose that
over a field extension $\Omega/k$ its diagonal factors in rational
Chow through finitely many products of smooth projective curves and
surfaces. Then such a factorization exists over a finite extension
$E/k$, with auxiliary curves and surfaces allowed to change.
\end{lemma}

\begin{proof}
Choose cycle representatives for the correspondences and finitely
many principal divisors witnessing the diagonal equality. All these
data, together with the auxiliary varieties and their projective
embeddings, descend to a field finitely generated over $k$. Spread
them over a smooth integral finite-type $k$-scheme $B$. Shrink $B$
so that the auxiliary families are smooth projective with
geometrically integral fibers, the cycles have their expected
codimensions, and the relative Chow equality holds. The last step
follows either by spreading its principal-divisor certificate or by
the direct-limit description of Chow groups over the function field.

Choose a closed point $b\in B$. Its residue field $E$ is finite over
$k$. Since $B$ is smooth, refined pullback to $b$ commutes with
the smooth pullbacks, proper projections and diagonal intersections
which define correspondence composition. The constant family
$X\times B$ specializes to $X_E$, while the auxiliary families
specialize to smooth curves and surfaces. Thus the diagonal
certificate specializes with all its degree shifts.
\end{proof}

\begin{theorem}[Complex certificates and surface moduli]\label{thm:moduli}
Let $X/F$ be smooth projective over a number field. If for one complex
embedding $h(X_\C)$ retracts through products of curves and surfaces
in rational Chow motives, then $X$ has a full height domain, also
after every finite number-field extension. In particular this holds
if $X_\C$ is a smooth projective moduli space appearing in
B\"ulles's theorem: Gieseker-stable twisted sheaves, or stable objects
for a generic stability condition, on a complex projective K3 or
abelian surface. Finite products of these spaces and the varieties
in \thmref{thm:hilbert} are allowed, with an everywhere potentially
good projective coefficient.
\end{theorem}

\begin{proof}
By \lemref{lem:specialize-certificate}, a finite extension $E/F$
carries an actual diagonal certificate of the required kind.
Complete it at a place above each finite place of $F$ and apply
\corref{cor:surface-envelope}; full-degree descent gives the
original local kernel vanishing. The same construction over a
compositum treats every new cycle after finite extension. Apply
\lemref{lem:two-degrees}.

For the moduli spaces, \cite[Theorem 0.1]{Bulles} supplies precisely
an actual rational Chow retraction through powers of the underlying
surface. The first assertion applies. For products, tensor the
certificates; the intermediate varieties remain products of curves
and surfaces. The surface's finite-dimensionality is not an
assumption of the cited retraction theorem.
\end{proof}



% Source: papers/pairing-sections/niveau_fibrations.tex
\section{Chow niveau, complete intersections, and flat fibrations}\label{sec:niveau}

\subsection{A residual motive of small dimension}

We use covariant rational Chow motives, with
$\CH_a(h(T)(u))=\CH_{a-u}(T)_\Q$.

\begin{theorem}[Chow-niveau transfer]\label{thm:niveau}
Let $X/k$ have dimension $n$. Suppose that after a finite extension
there is an actual rational Chow decomposition
\begin{equation}\label{eq:niveau-decomposition}
 h(X)\simeq Q\oplus(T,e,u),\qquad
 Q\text{ of abelian type},\quad \dim T\leq n-2u,
\end{equation}
where $u\geq0$. Over a finite extension of $\Q_p$ this gives zero
obstruction on the full geometric homological Chow groups for
\begin{equation}\label{eq:niveau-local}
 r\leq u+1\quad\text{or}\quad r\geq n-u-1.
\end{equation}
Every degree is covered if $n\leq2u+3$. Over a number field the full
height pairs are obtained in every degree in that case; otherwise
they are obtained for
\begin{equation}\label{eq:niveau-height}
 r\leq u+1\quad\text{or}\quad r\geq n-u.
\end{equation}

A sufficient hypothesis for \eqref{eq:niveau-decomposition} is a
morphism $N\to h(X)$ from an abelian-type Chow motive which, on an
algebraically closed universal domain $\Omega/k$ of infinite
transcendence degree over the field of definition, surjects on
algebraically trivial $\CH_i$ for every $i<u$.
\end{theorem}

\begin{proof}
Let $m=\dim T$. The residual image of a codimension-$r$ class has
dimension $n-r-u$ and codimension
\begin{equation}\label{eq:residual-degree}
 m-n+r+u\leq r-u.
\end{equation}
If $r\leq u+1$, this is a divisor degree or lower. If
$r\geq n-u-1$, the residual cycle dimension is at most one.
The endpoint input applies in either case. If $n\leq2u+3$, then
$m\leq3$, so every residual degree is covered. The abelian-type
summands are covered by \thmref{thm:main}. Their actual inclusion
and retraction maps, followed by \lemref{lem:general-transfer},
prove the local assertion. Intersect its range with the range of
$n+1-r$ to obtain \eqref{eq:niveau-height}.

For the sufficient hypothesis apply Vial's
\cite[Theorem 3.11]{Vial} with algebraic equivalence and $M=h(X)$.
The finite-dimensionality required there holds for $N$. Self-duality
of $h(X)$ gives \eqref{eq:niveau-decomposition} over the field of
definition of the morphism. Its additional curve and Tate summands
are of abelian type. A negative residual dimension means that the
residual summand is zero.
\end{proof}

In particular it suffices that
$\CH_i(X_\Omega)_{\Q,\alg}=0$ for $0\leq i<u$: take $N=0$.
This criterion concerns the subgroup algebraically equivalent to
zero, rather than a quotient by algebraic equivalence.

\subsection{Small-degree Chow bounds}

Let $X\subset\PP^N$ be a smooth complete intersection of degrees
$d_1\geq\cdots\geq d_c\geq2$, and put $n=N-c$.

\begin{proposition}\label{prop:small-degree}
Let $0\leq t\leq n-1$. Suppose either
\begin{equation}\label{eq:elv}
 \left(d_1\geq3\ \text{or}\ c\geq t+1\right),\qquad
 \sum_{a=1}^c\binom{d_a+t}{t+1}\leq N,
\end{equation}
or all $d_a=2$, $1\leq c\leq t$, and
\begin{equation}\label{eq:elv-quadratic}
 c(t+2)\leq N+c-t-1.
\end{equation}
Then \thmref{thm:niveau} applies with $u=t+1$.
For a hypersurface of degree $d\geq3$, an alternative sufficient
condition is
\begin{equation}\label{eq:hi-bound}
 (j+1)(N-j)+1\geq\binom{d+j+1}{j+1}
             \quad(0\leq j\leq t).
\end{equation}
In each case all degrees are covered if $n\leq2t+5$; otherwise
the guaranteed full height degrees are $r\leq t+2$ or
$r\geq n-t-1$.
\end{proposition}

\begin{proof}
Theorem 4.6 of Esnault--Levine--Viehweg \cite{ELV} gives
$\CH_i(X_\Omega)_\Q=\Q$ for $0\leq i\leq t$ under
\eqref{eq:elv} or \eqref{eq:elv-quadratic}. In the hypersurface
case, Hirschowitz--Iyer \cite[Theorem 1.5]{HI} gives the same
conclusion degree by degree under \eqref{eq:hi-bound}; the upper
variety is projective space and its relevant Chow group is
one-dimensional. A linear-section generator has nonzero projective
degree. The algebraically trivial subgroup is therefore zero in
all these degrees. Apply \thmref{thm:niveau} with $u=t+1$.
\end{proof}

There is a further useful bound for mixed quadrics and cubics.
Otwinowska's rewritten author version \cite[Corollaire 1]{Otwinowska}
gives the following statement for smooth complete intersections:
if there are $a$ quadrics and $b\geq1$ cubics in $\PP^N$, $N\leq15$,
then $\CH_i(-)_\Q$ is generated by the ambient section class for
\begin{equation}\label{eq:ot-bound}
 0\leq i\leq t=\left\lfloor\frac{N-2a-3b}{3}\right\rfloor.
\end{equation}
Its complex-field formulation gives the required universal-domain
statement in characteristic zero. Indeed, descend a given cycle
and its equations to a finitely generated field, embed that field
in $\C$, and obtain the rational-equivalence certificate there.
The certificate descends to a further finitely generated field,
which embeds over the original field of definition into the
universal domain. Its base change proves the original relation.
The coefficient of the ambient class is fixed by projective degree.

\begin{corollary}[Explicit complete-intersection domains]\label{cor:ci-list}
Every smooth geometrically integral variety in the following table
has $\mathsf V$ locally and a full height domain over a number
field. A presentation after a finite extension suffices.
\begin{center}
\begin{tabular}{lll}
\toprule
Type & Dimension & Low-Chow parameter $t$\\
\midrule
Cubic hypersurface & $6,7$ & $1$\\
Cubic hypersurface & $8,9$ & $2$\\
Cubic hypersurface & $11$ & $3$\\
Quartic hypersurface & $7$ & $1$\\
Complete intersection $(2,3)$ & $6,7$ & $1$\\
Complete intersection $(2,3)$ & $9$ & $2$\\
Complete intersection $(2,2,3)$ & $7$ & $1$\\
Complete intersection $(3,3)$ & $7$ & $1$\\
\bottomrule
\end{tabular}
\end{center}
The cubic cases in dimension at most five already follow from
\corref{cor:rc-fibrations}.
\end{corollary}

\begin{proof}
For cubics the values in \eqref{eq:hi-bound} are
$4,10,20,35$ for $j=0,1,2,3$. The requisite minimum ambient
dimensions for $t=1,2,3$ are $6,9,12$, respectively. These imply
the first three rows. For the quartic row, $N=8$ and the two
inequalities are $9\geq5$ and $15\geq15$.

For $(2,3)$ of dimension $6$ or $7$, the zero-cycle statement
follows from \eqref{eq:elv}. The upper quadric, which may be
singular, has $\CH_2=\Q$ by the all-intersection form of
\cite[Theorem 4.6]{ELV} and \eqref{eq:elv-quadratic}.
The distinct-largest-degree
case of \cite[Theorem 1.5]{HI}, with $j=1,c=2$, has
$\rho=3N-19$ and $\rho+1\geq N-2$ exactly when $N\geq8$.
Thus $\CH_1=\Q$. The last three rows follow from
\eqref{eq:ot-bound}: their ambient dimensions are $11,10,9$,
respectively. In every row $n\leq2t+5$.
\end{proof}

\subsection{The all-fiber criterion}

\begin{proposition}[Flat-fibration transfer]\label{prop:fibration}
Let $f:Y\to B$ be flat, projective and surjective, with $Y,B$
smooth projective and $\dim Y=n$, $\dim B=b$. Fix an integer
$u\geq1$ with $n-2u\geq0$. Suppose that for
every geometric scheme fiber over a universal domain
\[
 \CH_i(Y_\omega)_\Q=\Q\quad(0\leq i<u).
\]
Assume $\mathsf V(B)$ in the native local theory. Then the full
local source vanishes for $r\leq u+1$ or $r\geq n-u-1$, and in
all degrees if $n-2u\leq3$. If these base hypotheses hold at every
finite place of a number-field model, the corresponding full
height degrees are obtained by \eqref{eq:range-intersection}.
\end{proposition}

\begin{proof}
Vial's flat-fibration decomposition \cite[Theorem 4.2]{VialFibrations}
is an actual Chow decomposition
\begin{equation}\label{eq:fibration-motive}
 h(Y)\simeq\bigoplus_{j=0}^{n-b}h(B)(j)\oplus(T,e,u),
 \qquad\dim T=n-2u.
\end{equation}
The base summands are admitted by hypothesis. The residual-degree
calculation \eqref{eq:residual-degree} and the endpoint input give
the stated local degrees. Transfer through the actual diagonal
and apply \lemref{lem:two-degrees} globally.
\end{proof}

\begin{corollary}[Quadric and cubic fibers]\label{cor:quad-cubic-fibers}
The following flat projective surjective maps, with smooth projective
geometrically integral total space and base, have $\mathsf V$:
\begin{enumerate}[label=\textup{(\roman*)}]
\item Every geometric scheme fiber is a nonzero quadric of fixed
dimension $q\geq1$ in $\PP^{q+1}$, and $1\leq b\leq3$.
The same holds for $b=4$, $q$ odd, if the base has $\mathsf V$;
in particular the base may be $\PP^4$.
\item Every geometric scheme fiber is a nonzero cubic hypersurface
of fixed dimension $q$, and
\[
 (q,b)=(5,1),(5,2),(6,1),\text{ or }(8,1).
\]
\end{enumerate}
The conclusions concern the entire geometric homological Chow
groups. Over number fields the same geometric hypotheses give
full height domains, with the base condition imposed at all
finite places when $b=4$.
\end{corollary}

\begin{proof}
A nonzero quadric scheme of dimension $q$ over an algebraically
closed characteristic-zero field has $\CH_i\simeq\Z$ for
$i<q/2$. This includes rank-one quadrics: Chow groups are
unchanged by reduction. For other singular quadrics, the cone
description and the split-quadric cellular filtration give the
same range; the cone localization proof is recorded in
\cite{InternalFibrations}. Take $u=\lceil q/2\rceil$ in
\propref{prop:fibration}. Its residual dimension is $b$ for even
$q$ and $b-1$ for odd $q$, giving (i).

For cubics of dimension at least five, the all-intersection
ELV bound gives $\CH_0=\CH_1=\Q$. For $q=8$, the HI bound
also gives $\CH_2=\Q$. These cited bounds apply to the reduced
singular hypersurfaces in question. A nonreduced cubic has
reduction a hyperplane or a union of two hyperplanes; the closed
union Chow sequence gives the same low-degree groups. Thus the
four cases have residual dimensions $2,3,3,3$, respectively,
using $u=2,2,2,3$. The base has dimension at most two and hence
has $\mathsf V$. Apply \propref{prop:fibration}.
\end{proof}

\subsection{A signed retraction for intersections of quadrics}

\begin{theorem}[Intersections of quadrics]\label{thm:quadrics}
A smooth complete intersection of $c$ quadrics in $\PP^N$,
$1\leq c\leq N-1$, has $\mathsf V$ in all degrees when
\[
 c\leq4,\qquad\text{or}\qquad c=5\text{ and }N\text{ is even}.
\]
For arbitrary $c$, the local degrees supplied are
\begin{equation}\label{eq:quadrics-local}
 r\leq\lfloor N/2\rfloor-c+2
 \quad\text{or}\quad r\geq\lceil N/2\rceil-2.
\end{equation}
The full global height degrees outside the full case are
$r\leq\lfloor N/2\rfloor-c+2$ or
$r\geq\lceil N/2\rceil-1$.
\end{theorem}

\begin{proof}
For $c=1$, pass to a finite splitting field of the smooth quadric.
Its cellular Chow decomposition is Tate, so \thmref{thm:main}
gives every degree. Assume $c\geq2$. Write the equations as
$Q_1,\ldots,Q_c$ and let
\[
 Y=\{\textstyle\sum_jt_jQ_j=0\}\subset\PP^N\times\PP^{c-1}.
\]
It is a flat family of nonzero $(N-1)$-dimensional quadric
schemes. The total space is smooth: at a possible singular point
all $Q_j$ vanish, while their independent conormal differentials
on the smooth complete intersection make $\sum_jt_jdQ_j$
nonzero. Set $m=c-1$, $E=X\times\PP^m$, and write
$j:E\hookrightarrow Y$, $p:E\to X$. The normal sequence is
\[
 0\longrightarrow N_{E/Y}\longrightarrow
 p^*\mathcal O_X(2)^{\oplus c}\longrightarrow
 p^*\mathcal O_X(2)\otimes\mathcal O_{\PP^m}(1)
 \longrightarrow0.
\]
Its degree-$m$ Chern class has fiber-leading term
$(-1)^mc_1(\mathcal O_{\PP^m}(1))^m$, and fiber integration gives
\[
 p_*c_m(N_{E/Y})=(-1)^m.
\]
Consequently the correspondences
\begin{equation}\label{eq:quadrics-retraction}
 U=j_*p^*,\qquad V=(-1)^mp_*j^!,\qquad VU=1
\end{equation}
give a retraction, by self-intersection. The calculation after
every smooth projective product shows the same identity for
Chow correspondences. Here $U$ sends codimension $r$ to $r+m$.

Apply \propref{prop:fibration} to $Y\to\PP^{c-1}$, with
$u=\lceil(N-1)/2\rceil$. Its residual dimension is $c-2$ if
$N$ is even, and $c-1$ if $N$ is odd. These are at most three
in the asserted full cases. Substituting $r+m$ in the outer
range and using \eqref{eq:quadrics-retraction} gives
\eqref{eq:quadrics-local}. Intersecting the two complementary
ranges gives the height statement.
\end{proof}

\begin{theorem}[Relative quadrics]\label{thm:relative-quadrics}
Let $B$ be smooth projective geometrically integral of dimension
$b\geq1$, and let $\mathcal V,\mathcal F$ be vector bundles of
ranks $N+1,c$, with $c\geq1$ and $N\geq c+1$.
Let $\pi:\PP_B(\mathcal V)\to B$ parametrize lines. Suppose that
the zero scheme of a regular section of
$\mathcal O(2)\otimes\pi^*\mathcal F$ is smooth projective
geometrically integral, is flat over $B$, and has expected
codimension $c$ in every geometric fiber. If $\mathsf V(B)$,
then it has $\mathsf V$ whenever
\[
 b+c\leq5\quad(N\text{ even}),\qquad
 b+c\leq4\quad(N\text{ odd}).
\]
In general the local degree range is
\[
 r\leq\lfloor N/2\rfloor-c+2
 \quad\text{or}\quad r\geq b+\lceil N/2\rceil-2.
\]
Over a number field with the base admitted at every finite place,
the corresponding full height range ends instead at
$r\geq b+\lceil N/2\rceil-1$. The base condition is automatic
for $1\leq b\leq3$.
\end{theorem}

\begin{proof}
Form $C=\PP_B(\mathcal F^\vee)$ and the universal quadric
$Y\to C$. Expected fiber codimension ensures that no nonzero
coefficient vector gives the zero quadratic equation. Smoothness
of the total space follows from the absolute differential of the
regular section defining $X$; the base directions are included.
The projective-bundle formula gives $\mathsf V(C)$. For
$E=X\times_BC$, the same normal sequence and fiber integration
as in \eqref{eq:quadrics-retraction} give a signed retraction with
shift $c-1$. The residual dimension of the universal quadric is
$b+c-2$ for even $N$ and $b+c-1$ for odd $N$. Apply the
flat-fibration theorem, transfer through the retraction, and use
\lemref{lem:two-degrees}.
\end{proof}



\section{Actual Chow identities}\label{sec:chow}

For the algebraic arguments, let $k$ be any characteristic-zero
field. Write $\mathcal A_k(V)$ for a diagonal factorization of $V$
through abelian varieties over $k$ itself, in the sense of
\eqref{eq:diagonal-intro} with $L=k$. Smooth disconnected varieties
are interpreted by the direct sum of their components.

\begin{lemma}\label{lem:tensor}
The property $\mathcal A_k$ is stable under finite products, finite
disjoint unions of pure-dimensional components, isomorphisms and
base extension. A pointed smooth projective geometrically integral
curve has this property.
\end{lemma}

\begin{proof}
For products, take exterior products of the two sets of
correspondences, with the factors ordered as source and target.
Composition commutes with exterior products, so
\[
 \Delta_{V\times W}=\Delta_V\boxtimes\Delta_W
       =\sum_{j,l}(v_j\boxtimes v_l)\circ(u_j\boxtimes u_l).
\]
The intermediate variety $A_j\times A_l$ is abelian. This equality
is in the Chow group of the product with itself. Disjoint unions
and isomorphisms follow by inserting component maps and by conjugating
with graphs; flat field extension preserves the identity.

For a pointed curve $C$, use
$h(C)=\mathbf1\oplus h^1(C)\oplus\LL$ and the identification of
$h^1(C)$ with the corresponding summand of $h(J(C))$. These are
Chow-motive decompositions \cite[\S3.2, Proposition 4.5 and
Theorems 5.2--5.3]{Scholl}. Their inclusion and retraction maps,
with the Tate terms through $\Spec k$, give the required identity.
\end{proof}

\subsection{A codimension-two blowup}

Let $j:C\hookrightarrow V$ be a smooth center of codimension two,
and let
\[
 q:W=\Bl_CV\longrightarrow V,\qquad i:E\hookrightarrow W,
 \qquad p:E=\PP_C(N_{C/V})\longrightarrow C.
\]
Set $\xi=c_1(\mathcal O_E(1))$. Then $N_{E/W}=\mathcal O_E(-1)$
and $p_*\xi=1$.

\begin{lemma}\label{lem:blowup}
There is an equality of rational Chow correspondences
\begin{equation}\label{eq:blowup-diagonal}
        \Delta_W=q^*q_*-i_*p^*p_*i^*.
\end{equation}
If $\mathcal A_k(V)$ and $\mathcal A_k(C)$ hold, so does
$\mathcal A_k(W)$.
\end{lemma}

\begin{proof}
The Chow blowup formula \cite[\S6.7]{Fulton} writes every class
$z\in\CH^r(W)_\Q$ uniquely as
$z=q^*a+i_*p^*b$, where $b\in\CH^{r-1}(C)_\Q$. The projection
and self-intersection formulas give
\[
\begin{array}{ll}
 q_*q^*a=a,&q_*i_*p^*b=0,\\
 p_*i^*q^*a=0,&p_*i^*i_*p^*b=-p_*(\xi p^*b)=-b.
\end{array}
\]
Consequently $a=q_*z$ and $b=-p_*i^*z$, which proves
\eqref{eq:blowup-diagonal} as an operator identity.

To obtain a correspondence identity, repeat the argument after
product with any smooth projective $T$. Flat product commutes with
the blowup, and the same calculation applies to every Chow group
of $T\times W$. Taking $T=W$ and applying the operator on the
second factor to $\Delta_W$ proves the equality of correspondences.
Finally insert the given factorizations of $V$ and $C$. The center
maps $-p_*i^*:W\rightsquigarrow C$ and
$i_*p^*:C\rightsquigarrow W$ have degrees $-1$ and $1$.
\end{proof}

\subsection{Generically finite morphisms}

\begin{lemma}\label{lem:degree}
Let $g:W\to X$ be a dominant generically finite morphism of degree
$d$ between smooth projective geometrically integral varieties of
the same dimension. Then
\begin{equation}\label{eq:degree-diagonal}
                       g_*g^*=d\Delta_X
\end{equation}
as Chow correspondences. Thus $\mathcal A_k(W)$ implies
$\mathcal A_k(X)$.
\end{lemma}

\begin{proof}
For every smooth projective $T$, the Chow projection formula gives
$(g_T)_*g_T^*z=z\cdot(g_T)_*[T\times W]=dz$.
Taking $T=X$ and $z=\Delta_X$ proves \eqref{eq:degree-diagonal}.
If $\Delta_W=\sum_jv_ju_j$, then
\begin{equation}\label{eq:degree-maps}
             \Delta_X=\sum_j(d^{-1}g_*v_j)\circ(u_jg^*).
\end{equation}
The maps in each pair have opposite degrees.
\end{proof}

\section{Separated binary hypersurfaces}\label{sec:binary}

Allow first $2b+a\ge2$ in \eqref{eq:binary-family}, so that
zero-dimensional factors can occur in an induction. These factors
are finite \'etale. For every binary block, Euler's identity
shows that simultaneous vanishing of its derivatives at a nonzero
vector would give a multiple projective root. A nonzero unary
term has nonvanishing derivative away from zero in characteristic
zero. The full gradient therefore vanishes only at the affine
origin. The hypersurface is smooth; when its dimension is positive,
connectedness of a projective hypersurface makes it geometrically
integral.

\subsection{The controlled join}

Let $F(x_0,\ldots,x_r)$ and $G(y_0,\ldots,y_s)$ be homogeneous of
degree $d$, with $r,s\ge1$, such that $X_F$ and $X_G$ are smooth.
Set
\[
 A=\{F+u^d=0\},\quad B=\{G+v^d=0\},\quad H=\{F+G=0\}.
\]

\begin{lemma}\label{lem:join}
Over a field containing $\zeta$ with $\zeta^d=-1$, the rational map
\begin{equation}\label{eq:join-map}
 A\times B\dashrightarrow H,\qquad
 ((x,u),(y,v))\longmapsto(vx,\zeta uy)
\end{equation}
extends to a morphism
$g:\Bl_{X_F\times X_G}(A\times B)\to H$ of generic degree $d$.
The center is smooth of codimension two.
\end{lemma}

\begin{proof}
Substitution in the equation of $H$ gives
$v^dF(x)+\zeta^du^dG(y)=-u^dv^d+u^dv^d=0$.
Near a base point, choose nonzero coordinates $x_i,y_j$ and
trivialize the projective line bundles. The base ideal of the
displayed sections is then $(u,v)$. These are local equations
of transverse smooth divisors from the two factors. Their blowup
makes the base ideal invertible and resolves the map, with center
$X_F\times X_G$.

On the dense open of $H$ where $F(X)=-G(Y)\ne0$, normalize
$u=v=1$. A preimage has
$x=\lambda X$, $y=\lambda\zeta^{-1}Y$, and
$\lambda^d=-1/F(X)$. There are exactly $d$ distinct geometric
solutions. The exceptional divisor has smaller dimension and
does not dominate this open. Hence $g$ is generically finite
of degree $d$. This is the controlled construction of
\cite[Remark 1.10]{SK}, with its center and degree retained.
\end{proof}

\subsection{The recursive diagonal}\label{proof:binary}

We prove that every separated binary/unary hypersurface has
$\mathcal A$ after a finite extension. In dimension zero, split
the finite \'etale scheme. In dimension one, choose a
point after a finite extension and apply \lemref{lem:tensor}.

Suppose $H$ has dimension $n\ge2$. Choose one binary block as
$G$; if all blocks are unary, combine two unary terms into $G$.
The remaining polynomial $F$ has $n\ge2$ variables. In the
join construction, the dimensions are
\[
 \dim A=n-1,\quad \dim B=1,\quad
 \dim X_F=n-2,\quad \dim X_G=0.
\]
All four varieties belong to the same class. Induction, followed
by one common finite extension, supplies actual factorizations
for $A$, $B$, $X_F$ and every component of $X_G$. We include
the root needed for \lemref{lem:join} in that extension.

The tensor identity gives factorizations for $V=A\times B$ and
$C=X_F\times X_G$. \lemref{lem:blowup} gives one for $W=\Bl_CV$,
and \lemref{lem:degree} gives one for $H$. All branches of the
induction have strictly smaller single-factor dimension; finitely
many finite extensions suffice. Products and finite forms follow
from \lemref{lem:tensor}. This proves part (i) of
\thmref{thm:families}.

More explicitly, suppose
$\Delta_V=\sum_jv_j^Vu_j^V$ and $\Delta_C=\sum_l v_l^Cu_l^C$.
The maps for $H$ are
\begin{align}
 u_j^H&=u_j^Vq_*g^*,& v_j^H&=d^{-1}g_*q^*v_j^V,\label{eq:ambient-recursion}\\
 u_l^H&=-u_l^Cp_*i^*g^*,& v_l^H&=d^{-1}g_*i_*p^*v_l^C.
 \label{eq:center-recursion}
\end{align}
Substituting \eqref{eq:blowup-diagonal} and
\eqref{eq:degree-diagonal} verifies their sum is $\Delta_H$.

\begin{example}\label{ex:quartic-sixfold}
For $\lambda_i\in k\setminus\{2,-2\}$, the quartic sixfold
\[
 \sum_{i=0}^{3}
 \bigl(x_{2i}^4+\lambda_i x_{2i}^2x_{2i+1}^2+x_{2i+1}^4\bigr)=0
       \quad\text{in }\PP^7_k
\]
belongs to the family. A binary block has a repeated projective
root precisely when $\lambda_i=2$ or $-2$. Over a finite extension
of $\Q_2$, the same characteristic-zero calculation applies.
\end{example}

\section{Split cubics and finite Chow descent}\label{sec:cubic}

Laterveer proves that the cubic varieties
\eqref{eq:cubic-family} over $\C$ have Chow motives of abelian
type \cite[Theorem 5]{Laterveer}. This is a conclusion about
rational equivalence, and gives a retract of a finite sum of
Tate twists of abelian motives. We need its finite field of
definition, including the rational equivalences in the retraction.

\begin{lemma}\label{lem:descent}
Let $k_0/\Q$ be finitely generated, let $V_0/k_0$ be smooth
projective, and fix an embedding $k_0\hookrightarrow\C$.
If $h(V_{0,\C})$ is of abelian type, then an actual abelian
Chow retraction of $V_0$ is defined over a finite extension
$k_1/k_0$. For every embedding $k_0\hookrightarrow k$ into
a characteristic-zero field, it therefore exists over a finite
extension of $k$.
\end{lemma}

\begin{proof}
Vial's descent theorem \cite[Theorem 1 and \S1.1]{Vial} applies
to the extension $\bar k_0\subset\C$. It shows that
$h(V_{0,\bar k_0})$ is of abelian type. By the definition of
this category, it is a retract of a finite sum of Tate twists
of motives of abelian varieties: tensor products remain abelian
motives, and duality is absorbed by a Tate twist. Thus there
are finitely many correspondences satisfying the diagonal
identity over $\bar k_0$.

The varieties, group laws, cycles, and rational equivalences in
that identity are all defined by finitely many coefficients.
After clearing rational denominators, each rational equivalence
is a finite sum of divisors of rational functions on integral
subvarieties. These data descend to one finite subextension
$k_1/k_0$. Equivalently, the finitely many equalities in the
Chow groups hold at a finite stage of their directed limit over
$\bar k_0$.

Extend $k_0\hookrightarrow k$ to $k_1\hookrightarrow\bar k$.
Then $L=kk_1$ is finite over $k$, and flat base change carries
the same Chow identity to $L$.
\end{proof}

\subsection{Application to the block family}\label{proof:cubic}

For a cubic over $k$ as in \eqref{eq:cubic-family}, put its
coefficients in a finitely generated subfield $k_0\subset k$.
The smoothness assumptions descend and remain true after
embedding $k_0$ into $\C$. Laterveer's theorem applies there,
and \lemref{lem:descent} gives the required finite Chow
factorization over $k$. Tensoring the factorizations deals
with products; transporting them through an actual isomorphism
deals with finite forms. This proves part (ii) of
\thmref{thm:families}.

For example, two arbitrary smooth cubic threefold blocks followed
by a smooth binary cubic give a cubic tenfold in $\PP^{11}$.
Replacing the binary block by a smooth four-variable cubic
gives a cubic twelvefold in $\PP^{13}$. Their central cycle
codimensions are $5$ and $6$. The two five-variable blocks
need not themselves separate into binary terms.

\section{Separated-branch cyclic covers}\label{sec:cyclic}

Use the data of \eqref{eq:cyclic-family}, set $d=qm$, and put
\[
                  Y=\{T^d=F(x_0,\ldots,x_n)\}\subset\PP^{n+1}_k.
\]

\begin{lemma}\label{lem:cyclic-cover}
The varieties $X=X(F;q,m)$ and $Y$ are smooth, projective and
geometrically integral of dimension $n$. The map
\begin{equation}\label{eq:finite-cyclic}
                    f:Y\longrightarrow X,\qquad u=T^m,
\end{equation}
is finite flat of rank $m$.
\end{lemma}

\begin{proof}
The block derivative calculation in \secref{sec:binary} shows
that $B=\{F=0\}\subset\PP^n$ is smooth and nonempty. On the
chart $x_i\ne0$, write $R_i=k[x_j/x_i:j\ne i]$ and
$F_i=F/x_i^d$. The cover algebra is
\[
                      B_i=R_i[u_i]/(u_i^q-F_i).
\]
It is free of rank $q$ over $R_i$. Off $B$, the derivative in
$u_i$ is nonzero; above $B$, the differential of $F_i$ proves
smoothness. Over $\bar k$, the valuation at any component of
$B$ gives $F_i$ value one. Eisenstein's criterion shows that
$U^q-F_i$ is irreducible over the function field. The chart
algebras inject into this field, so the cover is geometrically
integral. Being finite over $\PP^n$, it is projective.

The polynomial $T^d-F$ has the separated-block form with one
additional unary term. Thus $Y$ is smooth and geometrically
integral. Its projection to $\PP^n$ has no base point. With
$t_i=T/x_i$, the source algebra is
\begin{equation}\label{eq:cyclic-algebra}
 R_i[t_i]/(t_i^{qm}-F_i)\simeq B_i[t_i]/(t_i^m-u_i).
\end{equation}
The right side is free of rank $m$ with basis
$1,t_i,\ldots,t_i^{m-1}$. The transition laws
$t_j=(x_i/x_j)t_i$ and $u_j=(x_i/x_j)^m u_i$ glue these maps.
This proves the rank everywhere, including the length-$m$
fibers with equation $t_i^m=0$ along the branch.
\end{proof}

\subsection{The product diagonal}\label{proof:cyclic}

Each $Y_i$ associated with a cover $X_i$ belongs to the separated
hypersurface family. Over one finite extension, choose its Chow
factorization $\Delta_{Y_i}=\sum_{j_i}v_{i,j_i}u_{i,j_i}$.
For a finite product put
$Y=\prod_iY_i$, $X=\prod_iX_i$, and $M=\prod_i m_i$.
The product map $f:Y\to X$ is finite flat of rank $M$.

For $J=(j_1,\ldots,j_s)$ set
\[
 A_J=\prod_i A_{i,j_i},\qquad
 u_J=\boxtimes_i u_{i,j_i},\qquad v_J=\boxtimes_i v_{i,j_i}.
\]
The tensor identity gives $\Delta_Y=\sum_Jv_Ju_J$.
Since $f_*f^*=M\Delta_X$, the maps
\begin{equation}\label{eq:cyclic-diagonal}
                   U_J=u_Jf^*,\qquad V_J=M^{-1}f_*v_J
\end{equation}
satisfy $\sum_JV_JU_J=\Delta_X$. The rank $M$ is included
once, and all exceptional signs from the factorization of $Y$
are retained. Actual finite forms are handled by transport of
this identity. This proves part (iii) of \thmref{thm:families}.


% Source: papers/pairing-sections/models_families.tex
\section{Model criteria and arithmetic families}\label{sec:models}

\subsection{Numerical corrections on a semistable witness}

The numerical-model criterion of \propref{prop:numerical} is often
applied through a projective dominant witness. In the general
companion formulation \cite{InternalEndpoints}, the witness carries
a rational virtual bundle whose generic degree-$r$ Chern character
is the pulled-back class and whose degree on every vertical
$r$-dimensional subvariety is zero. The native model-Chern detector
then gives zero obstruction upstairs, and dominant detection gives
zero obstruction downstairs. This formulation permits a witness
of positive relative dimension as well as a generically finite one.

We record three ways of producing this numerical correction.
The model calculations and their supporting residue sequences are
given in \cite{InternalModelCriteria}; all model conditions below
refer to supplied projective models.

\begin{proposition}[Scalar total boundary image]\label{prop:scalar-model}
Let $R=K^\circ$ and let $\mathcal X/R$ be integral, regular,
projective reduced strictly semistable, with smooth geometrically
integral generic fiber,
of relative dimension $n\geq3$, with smooth geometrically integral
components $D_i$, smooth closed strata, connected dual graph, and
one relatively ample class $H$. Fix $1\leq r\leq n$ and an
invertible prime $\ell$. Assume that at every vertex $i$ the sum
of Gysin images
\[
 \sum_{j\ne i}\im\bigl(
 H^{2r-2}(E_{ij,\bar k},\Q_\ell(r-1))
       \longrightarrow H^{2r}(D_{i,\bar k},\Q_\ell(r))\bigr)
\]
has dimension one. Then every geometrically homologically trivial
codimension-$r$ cycle on the generic fiber has zero obstruction.
\end{proposition}

\begin{proof}
Choose a Chow extension $\Theta$. The companion specialization
calculation puts its component class in the displayed total image.
Set
\[
 \lambda_i=\deg_k(i_i^!\Theta\,H_i^{n-r}),\qquad
 w_{ij}=\deg_k((H|_{E_{ij}})^{n-1}),
\]
summing disconnected intersections with their native residue degrees.
The specialization trace gives $\sum_i\lambda_i=0$.
The connected weighted graph has positive weights and hence a rational
solution of
\[
 \sum_jw_{ij}(v_j-v_i)=\lambda_i.
\]
The corrected extension
$\Theta_{\rm ad}=\Theta-\sum_i v_i i_{i*}H_i^{r-1}$ has zero
component class: in a one-dimensional total image its degree against
$H_i^{n-r}$ detects that class. Its pure operational representative
is numerically zero on every vertical $r$-cycle. Apply the numerical
model criterion. A constant change of all $v_i$ adds the principal
total-fiber class multiplied by $H^{r-1}$ and hence changes nothing
in Chow.
\end{proof}

One geometric instance is a supplied two-factor hypersurface
degeneration
\begin{equation}\label{eq:two-factor-model}
 AB+\varpi G=0\quad\text{in }\PP^{n+1}_R,
 \qquad\deg A=a,\quad\deg B=b,\quad\deg G=a+b.
\end{equation}
Assume the two special hypersurfaces, their intersection, and
$\Sigma=(A=B=G=0)$ have the smooth transverse geometrically
integral strata specified by these equations, and that the generic
fiber is smooth geometrically integral. Blowing up $(A,\varpi)$
gives components $\Bl_\Sigma(A=0)$ and $(B=0)$ meeting in
$(A=B=0)$. The blowup and restriction computation in
\cite{InternalModelCriteria} verifies the scalar-image condition in
every degree. Thus these local hypersurfaces have $\mathsf V$.

\begin{proposition}[Positive primitive components]\label{prop:positive-model}
Let a smooth projective fourfold have, after a finite extension, a
projective dominant four-dimensional witness with a regular flat
projective model. Write its special fiber as $\sum_i m_iD_i$, with
$m_i>0$, reduced strict normal-crossings support and smooth closed
strata. Suppose a relatively ample $H$ makes
\[
 \deg(xy)\quad\text{positive definite on}\quad
 P_H^2(D_i)=\ker\bigl(H:N^2(D_i)_\Q\to N^3(D_i)_\Q\bigr)
\]
for every component. Then the original fourfold has $\mathsf V$.
\end{proposition}

\begin{proof}
The companion numerical correction leaves a canonical primitive
residual $R_i$ with
$\sum_i m_i\deg(R_i^2)=0$. This residual identity, including the
native constant-field factors, is the primitive-residual theorem of
\cite{InternalModelCriteria}. Positivity and $m_i>0$ give $R_i=0$
for every $i$. The resulting operational lift is numerically
admissible, so codimension two follows by witness descent.
One-cycles and the endpoint degrees complete the fourfold.
\end{proof}

The verified component classes include abelian fourfolds, rank-two
projective bundles over smooth threefolds, rank-three bundles over
smooth surfaces, rank-four bundles over smooth curves, and $\PP^4$.
For the abelian case the positivity input is \cite{Ancona}; the
projective-bundle cases follow from the numerical projective-bundle
formula and divisor Hodge index. If $f:D'\to D_{k'}$ is finite
surjective and $(D',f^*H_{k'})$ has the required positivity, then
$(D,H)$ has it by pullback and the degree formula.

\begin{proposition}[Finite model symmetries]\label{prop:symmetry-model}
For a regular projective model of relative dimension $n\geq4$,
with the preceding multiplicity and stratum conventions, put
\[
 P_H^2=\bigoplus_i\ker\bigl(H^{n-3}:N^2(D_i)_\Q\to
                                    N^{n-1}(D_i)_\Q\bigr),
 \qquad Q(x,y)=\sum_i m_i\deg(x_iy_iH_i^{n-4}).
\]
Let a finite group act on the model preserving the numerical class
of $H$. If $Q$ is nondegenerate and $P_H^2$ is multiplicity-free
as a rational representation, every codimension-two homological
class has zero obstruction. In dimension four, nondegeneracy is
automatic and the conclusion gives $\mathsf V$.
\end{proposition}

\begin{proof}
The image of the canonical residual map is an invariant totally
isotropic submodule of $P_H^2$, by the equivariant residual theorem
in \cite{InternalModelCriteria}. Every rational simple representation
of a finite group is self-dual, by an averaged positive rational
form. Distinct simple summands are therefore orthogonal for $Q$.
Multiplicity-freeness and nondegeneracy imply that each summand is
nondegenerate, so no nonzero invariant totally isotropic submodule
exists. The residual vanishes. Numerical correction and witness
descent give the asserted degree. In dimension four, numerical
duality and divisor Hodge index give the required nondegeneracy;
the endpoints and one-cycles give all remaining degrees.
\end{proof}

The same proof applies to a sum of rational central isotypic
projectors $e$ when $eP_H^2$ has these properties, and gives the
corresponding cycle subspace $e\CH^2_{\Q,\geomhom}$. Such a
subspace pairs globally with the admitted complementary group.

\subsection{Ordinary quadratic reduction}\label{sub:ordinary-reduction}

\begin{theorem}[Companion ordinary-reduction theorem]\label{thm:ordinary}
Let $K/\Q_p$ be finite for any prime $p$, and let $X/K$ be smooth
projective geometrically integral of dimension at least three.
Suppose $X$ has an integral flat projective model whose geometric
special fiber is integral and smooth apart from finitely many
ordinary quadratic points, meaning that their projectivized
geometric tangent cones are smooth quadrics. Then $\mathsf V(X)$ holds. The same
conclusion holds if such a smooth projective dominant witness is
supplied after a finite extension. The theorem reapplies to all
new cycles over every finite valued extension.
\end{theorem}

\begin{proof}
In odd residue characteristic the normal-form and successive
reduced-point blowup calculation of \cite{InternalOrdinary}
produces a regular reduced quadric forest after finite extension.
Its numerical corrections are compatible with the incidence maps.
The all-characteristic quadric-forest theorem in the same source
handles odd-dimensional ordinary quadratic models at residue two.

For even $n\geq4$ at residue two, use the henselian normal form of
\cite[Lemma 4(ii)]{JannsenSaito}. A common finite extension splits
all the finitely many node polynomials. Translating each by an
integral root gives an actual pointed \'etale local equation
\[
 P(x_1,\ldots,x_n)+z^2+\epsilon\varpi^qz=0,
 \qquad q\geq1,\quad\epsilon\in R^\times,
\]
with unimodular polar form for $P$, where $R$ is the auxiliary
field's valuation ring and $\varpi$ its uniformizer. The common \'etale charts and
exact saturated blowup calculations are proved in
\cite[common \'etale neighborhoods and native degenerate-node chains]
{InternalOrdinary}. Each reduced-point blowup decreases $q$ by
one. After exactly $q$ blowups the exceptional chain consists of
$q-1$ copies of
$\PP_E(\mathcal O_E\oplus\mathcal O_E(-1))$ and a terminal
smooth quadric, where
\[
 E=(\overline P(X)+Z^2=0)\subset\PP^n,
 \qquad
 T=(\overline P(X)+Z^2+\overline\epsilon ZS=0)
       \subset\PP^{n+1}.
\]
All double strata are the smooth odd-dimensional quadric $E$;
all multiplicities are one and the intermediate sections are
disjoint. The same quadric-forest numerical criterion now applies.

In every case numerical witness descent uses the complete common
field degree, and the original curved pair gives the form on the
specified support complement. The dominant-witness version follows
from the same detector. After a finite valued extension the ordinary
geometric special points persist, with possibly larger finite
orders; apply the construction anew to every new cycle.
\end{proof}

Over a number field, good reduction or the ordinary alternative at
every finite completion, also on the stated dominant witnesses,
gives the full height pairing. Local auxiliary extensions are
descended before this global construction and need not come from
one fixed global extension.

\subsection{Tate and abelian-type strata}

\begin{theorem}[Companion all-strata criterion]\label{thm:abelian-strata}
Let $Y,T/K$ be smooth projective geometrically integral. Suppose
$Y$ has a supplied regular projective reduced strictly
semistable model whose every connected closed geometric stratum
has a rational Chow motive of abelian type after finite constants.
Assume weight--monodromy for $Y$ in every degree and every
invertible coefficient prime. If $T$ has potentially good projective
reduction, then $T\times Y$ has $\mathsf V$, also for all new
cycles after every finite valued extension.
\end{theorem}

\begin{proof}
The ramified all-strata theorem in \cite{InternalUniformization}
preserves the stated Chow-motive condition on a projective
semistable repair. Its toroidal calculation expresses new strata
through blowups of old boundary strata and projective bundles over
them. For a closed-boundary blowup the property holds in both
directions; this permits the toroidal factorization steps. The
character lattice of a new stratum is primitive, and its calculation
retains the actual line-bundle transition functions.

For a new cycle over $L/K$, pass to a finite extension where $T$
is good and the strata split. On one abelian carrier, numerical
semisimplicity and its comparison prime supply an algebraic
cohomological inner inverse to the boundary incidence. Tensor it
with the diagonal of the coefficient, so the inverse acts on all
mixed K\"unneth summands. Weight--monodromy gives the odd-weight
condition, while the inverse gives numerical-kernel lifting and
actual-span saturation. Apply the companion model criterion and
descend with the full degree. The final curved-pair correction is
on the original support of the newly specified cycle.
\end{proof}

\begin{corollary}[Drinfeld uniformization]\label{cor:drinfeld}
Let $X_a$ be smooth projective algebraizations of
$\Gamma_a\backslash\Omega_K^{d_a}$, where
$\Gamma_a\subset\operatorname{PGL}_{d_a+1}(K)$ is torsion-free,
discrete and cocompact, and $d_a\geq1$. Then
$T\times\prod_aX_a$ has $\mathsf V$ for every potentially good
projective coefficient $T$, also over every finite extension.
\end{corollary}

\begin{proof}
The projective semistable cover of a Drinfeld quotient has
components and closed strata obtained by the successive
rational-linear-subspace blowups of projective space. Their actual
Chow motives are split Tate, by the blowup formula. Staircase
products retain this property. Ito's theorem \cite{Ito} gives
weight--monodromy, and K\"unneth gives it for the product. Apply
\thmref{thm:abelian-strata} and descend the finite cover.
\end{proof}

One may also include curve or surface factors supplied with
projective strict semistable models all of whose closed strata are
of Chow abelian type. The resulting placewise presentations give
full global heights when available at every finite completion;
the factors in those local presentations may depend on the place.

\begin{theorem}[Totally degenerate strata]\label{thm:totally-degenerate}
Suppose a smooth projective $X/K$ has, over a finite extension
$L/K$, a projective dominant generically finite morphism
$Y\to X_L$, with $Y$ smooth projective geometrically integral
of dimension $\dim X$ and carrying a regular projective reduced
strictly semistable model. On every
connected geometric closed stratum $D$, require:
\begin{enumerate}[label=\textup{(\roman*)}]
\item Each integral Chow group is finitely generated. The compatible
hyperplane class $H$ from a projective embedding of the special
fiber satisfies, for $2c\leq\dim D$ and nonzero rational primitive
$x$ with $H^{\dim D-2c+1}x=0$,
\[
 (-1)^c\deg(xH^{\dim D-2c}x)>0.
\]
\item For every $\ell$ different from the residue characteristic,
odd geometric integral cohomology is torsion and
\[
 (\CH^c(D)/\mathrm{tors})\otimes\Z_\ell
 \simeq H^{2c}(D,\Z_\ell(c))/\mathrm{tors}.
\]
\end{enumerate}
Then $\mathsf V(X)$ holds. Smooth projective toric strata satisfy
these hypotheses. A global full height domain follows if such
witnesses, or another full local criterion, are supplied at every
finite place.
\end{theorem}

\begin{proof}
These are conditions (a),(b) of Raskind--Xarles
\cite[Definition 1]{RX}. Define finitely many Chow generators and
components over one finite residue extension and lift it
unramifiedly. The cycle isomorphisms and Poincar\'e duality identify
native numerical groups with even cohomology, compatibly with the
actual incidence correspondences. Thus every numerical-kernel
test lifts to a closed cohomological support test, by the companion
component-detection theorem \cite{InternalDegenerate}.

The cited Raskind--Xarles weight result kills the relevant odd
local Galois groups. More precisely, for
$W=H^{2t-1}(Y_{\bar L},\Q_\ell(t))$ the inertia invariants have
weights at most $-1$, and the inertia-$H^1$ term has weights at
least $1$. Hence $W^{G_L}=H^1(L,W)=0$. The companion period
comparison identifies the rational detector period with the
local Abel--Jacobi cup-product against the closed support test.
It vanishes in $\Q_\ell$, hence is zero as its original rational
intersection number. The finite detector and exact-support
reconstruction give the result. The global assertion applies
\lemref{lem:two-degrees} to these supplied local witnesses.
\end{proof}

\subsection{A common arithmetic construction}\label{sub:arithmetic-construction}

The dense families below use actual local models, together with an
arithmetic pencil argument. We give the latter separately to specify
the field and density quantifiers.

\begin{lemma}[Congruence pencils]\label{lem:congruence-pencils}
Let a fixed integral flat projective ambient model with geometrically
integral fibers and a section system be supplied over a number
ring. Suppose a finite free coefficient module surjects
onto the section system and its required residue spaces. Suppose
the system has a nonempty fixed-member Lefschetz-pencil open, and
that one can prescribe smooth sections at a finite set of places
and an established local smoothing model at distinguished places.
Require the chosen equations to be nonzero on every special
ambient fiber, and the generic smooth-member open to consist of
geometrically integral members.
After a single finite extension realizing these finite residue data,
the congruence-pencil construction gives a Zariski-dense set of
smooth generic members with the prescribed local models, smooth
reduction at the exceptional places, and smooth or ordinary-quadratic
reduction at all remaining places.
\end{lemma}

\begin{proof}
Choose a smooth reference section $H_0$ and include in a finite
set all rational primes below its bad reductions, its zero special
sections, and residue characteristic two. Prepare the desired
residue sections after finite residue extensions. A monic
polynomial with separable irreducible reductions of one common
degree realizes all these extensions in one number field, unramified
at the chosen primes. This field is fixed from now on.

In the free coefficient module impose the finite congruences for a
second section $H_1$, using the square of each native maximal
ideal to retain a prescribed first-order smoothing. For a
distinguished rational prime $p$ unramified in the starting
number field, this is the modulus $p^2$ used in the applications.
Intersect this
affine congruence lattice with the fixed-$H_0$ Lefschetz open. The
intersection is Zariski dense. For a chosen $H_1$, spread its pencil
outside a finite further set of primes. At those additional primes,
choose the parameter to specialize to the already smooth reference
member; at the originally prepared primes choose it to preserve
the chosen residue member. Chinese remainders give infinitely many
parameters with these requirements. Omit the finitely many singular
generic members.

Outside the spreading set, the pencil has smooth or ordinary
quadratic fibers. At every exceptional prime the selected residue
fiber is smooth or is the prescribed established smoothing.
The nonzero residue equations are non-zero-divisors on the integral
ambient fibers. The relative Cartier criterion therefore makes the
resulting divisors flat over the number ring. Generic geometric
integrality follows from the prescribed open. In each application,
nonzero residue equations are checked at the exceptional primes by
the congruences and elsewhere by independence of the two pencil
sections. Finally a polynomial
vanishing on all constructed members vanishes on each permitted
pencil, since that pencil contains infinitely many of them. Varying
$H_1$ in the dense congruence open proves density in the stated
section space. The finite spreading set may depend on $H_1$; the
fixed field does not.
\end{proof}

\subsection{Hypersurfaces and complete intersections}

The following are the geometric conclusions of the local model
calculations in \cite{InternalDense}. Their arithmetic proofs use
\lemref{lem:congruence-pencils}. Each constructed member has full
heights on all new homological classes after every finite extension
of its stated field.

\begin{theorem}[Dense hypersurface families]\label{thm:dense-hypersurfaces}
For every $n\geq3$ and $D\geq3$, there is a Zariski-dense subset
of the full degree-$D$ hypersurface parameter space over $\Q$
consisting of smooth geometrically integral varieties with full
height domains over every number field.

For $n\geq4$, fix positive $d_1,\ldots,d_m$, $m\geq3$,
$D=\sum_i d_i$, and suppose
\[
 b_n=\sum_{|I|=n+1}\prod_{i\in I}d_i-\binom{m-1}{n}>0.
\]
There is such a dense subset with a distinguished product-profile
degeneration of extreme monodromy rank $\rk N^n=b_n$. For linear
factors this gives $D\geq n+2$ and
$\rk N^n=\binom{D-1}{n+1}$.
\end{theorem}

\begin{proof}
For the ordinary family take
\[
 F_0=\sum_{j=0}^{n+1}X_j^D,\qquad
 J_D=\sum_{j=0}^n X_j^{D-1}X_{j+1}+X_{n+1}^D.
\]
At a prime dividing $D$, descending through the partial derivatives
shows $J_D$ is smooth. At every other prime $F_0$ is smooth.
Choose an integral $F_1$ congruent to $J_D$ at primes dividing
$D$ and in the fixed-$F_0$ Lefschetz open. The members
$DF_0+aF_1$, with Chinese-remainder conditions on $a$, have smooth
reduction at the finitely many exceptional primes, including two,
and ordinary reduction elsewhere. Apply \thmref{thm:ordinary}.
Varying $F_1$ proves the asserted full-space density.

For the product profile, prepare transverse integral factors
$P=\prod_i f_i$ and a smoothing section $Q$. Choose
$q>D$ outside their finite bad set and impose
$F_1\equiv P-DF_0+qQ\pmod{q^2}$ as well. For $a\equiv1$
modulo $q^2$, the exact equation at $q$ is $P+qQ_a$ with
$\bar Q_a=\bar Q$. The ordered product-model calculation in
\cite{InternalDense} gives full local vanishing and the displayed
extreme monodromy rank. Its ramified-extension theorem treats all
new cycles after arbitrary finite extension. At the other primes
use the preceding ordinary argument. The complementary-degree
criterion gives the heights.
\end{proof}

\begin{theorem}[Successive complete-intersection sections]\label{thm:dense-ci}
Let $\mathcal Y\subset\PP^N_{\mathcal O_F}$ be a supplied
integral flat projective relative complete intersection of dimension
$m\geq4$, with smooth geometrically integral generic fiber and
geometrically integral special fibers smooth outside finite sets.
For every $d\geq3$, after one fixed finite extension $E/F$ there
is a dense subset of the full degree-$d$ section system on $Y_E$
whose smooth members have full height domains over all finite
extensions of $E$.

Consequently one can construct complete intersections of any
specified multidegree $(d_1,\ldots,d_t)$ with all $d_i\geq3$
and final dimension at least three, over one fixed finite extension
of any prescribed number field. The first equations may be fixed
so that density holds in the final section system. A split quadric
may be used as the initial ambient variety.
\end{theorem}

\begin{proof}
The Koszul resolution supplies the full section space from ambient
polynomials. Smooth residue sections can be chosen to avoid the
finite ambient singular loci, after finite residue preparation.
The degree-$d$ jets for $d\geq3$ give the fixed-member Lefschetz
open. Apply \lemref{lem:congruence-pencils}; the resulting integral
complete-intersection model has smooth or odd-characteristic
ordinary fibers and good fibers at two. Use
\thmref{thm:ordinary}. Its model has the same finite-singular-locus
property, so one can iterate, fixing the earlier equations before
varying the last one.
\end{proof}

There is also a product-profile final-section version on the fixed
ambient of \thmref{thm:dense-ci}, with $m=n+1\geq5$. If the
ambient degree is $\delta$, the final degree is
$D=\sum_{i=1}^s e_i$ with $s\geq3$, $e_i\geq1$, the companion
construction gives a dense final-section family after one fixed
extension, with
\[
 \rk N^n=\delta\sum_{|I|=n+1}\prod_{i\in I}e_i
                                      -\binom{s-1}{n}.
\]
For example degree-five sections of the split quadric fivefold
give a dense family of type-$(2,5)$ fourfolds with $\rk N^4=1$.
The local product theorem and the same congruence-pencil proof
give their complete height domains.

\subsection{Product ambients and positive fourfold hosts}

\begin{theorem}[Dense product-ambient families]\label{thm:dense-products}
For each fixed set of input data below, the constructed dense set
has a full height domain over all finite extensions of one fixed
finite extension of its stated initial number field.
\begin{enumerate}[label=\textup{(\roman*)}]
\item For every $D\geq5$ and every prescribed number field
$F_{\rm base}$, one fixed finite $F/F_{\rm base}$ supports a
dense subset of the full bidegree-$(D,D)$ system on
$(\PP^1\times\PP^4)_F$. A distinguished product model has
\[
 \rk N^4=(D-1)\binom{D-1}{4}>0.
\]
\item Let $\mathcal S/\mathcal O_F$ be a supplied smooth
projective geometrically integral surface model and let $A$ be
relatively very ample. On $\mathcal S\times\PP^3$, put
$L=A\boxtimes\mathcal O(1)$. After one fixed finite extension
$E/F$, the image of degree-five embedding polynomials has a
dense such subset, with
$\rk N^4=10\deg(A^2)-1$. If those polynomials span the full
section space, this is density in $|L^5|$.
\item Let $\mathcal M/\mathcal O_{F_0}$ be a supplied smooth
projective scheme of relative dimension four over the whole number
ring, with geometrically integral fibers. Fix an odd rational prime $p$ unramified
in $F_0$. At each $v\mid p$ and every finite residue extension
$k'/k(v)$ assume
\begin{equation}\label{eq:host-signature}
 b^-(N^2(\mathcal M_{k'})_\R,\deg xy)
                  =\rho(\mathcal M_{k'})-1.
\end{equation}
Choose $H$ sufficiently positive for relative very ampleness,
normal generation and section base change. Put
$L_0=H\boxtimes\mathcal O_{\PP^1}(1)$ and $c=\deg H^4$.
One fixed finite $F/F_0$, unramified above $p$, supports a dense
subset of the full $|L_0^5|_F$ with
\[
 \rk N^4=4(5c-4)>0,\qquad
 b_1(X)=b_1(M),\quad b_2(X)=b_2(M)+1.
\]
\end{enumerate}
\end{theorem}

\begin{proof}
For (i), use $D$ transverse $\mathcal O(1,1)$ factors at a
distinguished prime. The ambient double-stratum comparison retains
both hyperplane classes. The ordered model and its uniform
finite-extension comparison, proved in \cite{InternalDense}, give
full local vanishing and
$5\binom D5-\binom{D-1}4=(D-1)\binom{D-1}4$.

For (ii), the local surface-ambient comparison uses the actual
blowups of the double strata. Its extreme incidence complex gives
$10\deg(A^2)-1$. The free module of embedding polynomials supplies
the coefficient space in \lemref{lem:congruence-pencils}, proving
density in its image. The stated spanning hypothesis identifies
that image with the full section system.

For (iii), the positive-host suspension calculation in
\cite{InternalPositiveHosts} applies to five divisors on $M$
and five distinct point factors on $\PP^1$. Residue Bertini
prepares smooth geometrically integral intersections of one,
two and three divisors, split finite-\'etale fourfold intersections
of degree $c$, and empty fivefold intersection. The smoothing
section cuts each positive-dimensional stratum transversely,
with geometrically integral positive-dimensional cuts and
finite-\'etale zero-dimensional cuts, and avoids the
zero-dimensional strata. Condition \eqref{eq:host-signature} kills
the primitive numerical residual. Its uniform form over finite
residue extensions gives the all-new-cycle ramified conclusion.
The incidence complex has $\dim H_3=5c-4$ before joining the
five point vertices, hence $\rk N^4=4(5c-4)$ afterward.
Weak Lefschetz gives the two Betti numbers. Relative Serre
vanishing and a sufficiently high Veronese supply the full
coefficient module and the required jets. Apply
\lemref{lem:congruence-pencils} to all three cases, use ordinary
reduction at nondistinguished bad primes, and then
\lemref{lem:two-degrees}.
\end{proof}

The signature hypothesis in (iii) is supplied by everywhere-good
abelian fourfolds, by the projective-bundle host types listed after
\propref{prop:positive-model}, and by $\PP^4$. For $\PP^4$,
$H=\mathcal O(1)$ and $p=5$, the result gives bidegree-$(5,5)$
fourfolds with $\rk N^4=4$, $b_1=0$, $b_2=2$.

\subsection{Explicit elliptic-power families}

Let $E_0/\Q$ be $y^2=x^3-x$ and put
$F_{\rm tor}=\Q(E_0[3])$, adjoining all geometric $3$-torsion.
The elliptic curve has an everywhere-good model $\mathcal E$
over this field. At two, potential good reduction follows from
complex multiplication; after adjoining $3$-torsion its finite
inertia image lies in the torsion-free principal congruence
subgroup of $\operatorname{GL}_2(\Z_3)$, so it is trivial. The
curve is already good at the other primes. This field is
unramified above five.

Put $A=\mathcal O_E(3[0])\boxtimes\mathcal O_E(3[0])$ on $E^2$,
and $H=\boxtimes_{i=1}^4\mathcal O_E(3[0])$ on $E^4$.
The plane-cubic embeddings and the Segre construction prove full
section spanning. The following dense subsets of full linear
systems are constructed over fixed finite extensions of
$F_{\rm tor}$ unramified above five, and have full heights after
every further finite extension:
\begin{center}
\begin{tabular}{llll}
\toprule
Ambient and system & Distinguished profile & Monodromy & $(b_1,b_2)$\\
\midrule
$E^2\times\PP^3$, $|(A\boxtimes\mathcal O(1))^5|$
 & five section factors & $\rk N^4=179$ & $(4,7)$\\
same system & curve and five hyperplanes & $N^4=0$ & $(4,7)$\\
$E^4\times\PP^1$, $|(H\boxtimes\mathcal O(1))^5|$
 & two curves and five points & $\rk N^2=802816$, $N^3=0$ & $(8,29)$\\
same system & five divisors and five points & $\rk N^4=38864$ & $(8,29)$\\
\bottomrule
\end{tabular}
\end{center}
The first row has $\deg A^2=18$. The second uses the mixed
surface--hyperplane local theorem of \cite{InternalMixedFamilies};
its weight-two contribution contains a nonzero eight-dimensional
summand. In the third row the double-stratum model is the actual
blowup of the curve product along the smoothing cuts. The local
signature condition holds for the abelian fourfold product, and
the exact rank formula is
\[
 \rk N^2=(2g(C)-b_1(S))(2g(D)-b_1(T))(d-1),\qquad N^3=0.
\]
Here $g(C)=g(D)=226$, $b_1(S)=b_1(T)=4$, and $d=5$.
The final row is \thmref{thm:dense-products}(iii), with
$c=4!\,3^4=1944$. The arithmetic argument is again
\lemref{lem:congruence-pencils}, now with the full Segre coefficient
system. The different monodromy profiles retain the distinctions
among these constructed subsets even when the ambient system is
the same.



% Source: papers/pairing-sections/cycle_domains.tex
\section{Cycle subgroups and explicit representatives}\label{sec:cycle-domains}

\subsection{Algebraically trivial cycles}\label{sub:algebraic-triviality}

Write $\CH^r(X)_{\Q,\alg}$ for the rational Chow classes whose
base change to an algebraic closure is algebraically equivalent
to zero.

\begin{proposition}\label{prop:algebraic-triviality}
On every smooth projective $X/K$ and in every codimension,
\[
 \CH^r(X)_{\Q,\alg}\subseteq\cK_K^r(X).
\]
Over a number field, the companion normalized height is therefore
defined on
\[
 \CH^r(X)_{\Q,\alg}\times
 \CH^{n+1-r}(X)_{\Q,\alg},
\]
and more generally on either algebraically trivial factor paired
with the admitted complementary group $\cH$.
\end{proposition}

\begin{proof}
Algebraic equivalence is generated by connected smooth projective
curve families. After a finite extension, an algebraically trivial
class is a finite sum of
\[
 \Gamma_*([a]-[b]),\qquad
 \Gamma\in\CH^r(C\times X)_\Q,
\]
with $a,b$ rational over that finite extension on the smooth
projective parameter curve. Enlarge the extension to retain the
finite rational-equivalence witnesses as well.
The degree-zero divisor $[a]-[b]$ is in its endpoint kernel.
Kernel covariance and full-degree finite descent prove the local
inclusion. The constant-field interpretation of algebraic
triviality is equivalently obtained by norming the finite family
certificate. Apply this inclusion at every completion and then
the companion height theorem.
\end{proof}

Together with the endpoint results, this also gives the full
divisor--zero-cycle height on an arbitrary $n$-fold. The complete
one-cycle group lies in $\cH^{n-1}(X)$ and pairs with
$\cH^2(X)$. Whenever one of the geometric or model criteria above
establishes the entire codimension-two group, this becomes the
full complementary pairing.

\subsection{Finite obstruction certificates}\label{sub:finite-certificates}

The finite-detector theorem of \cite{InternalDetectors} provides,
for a fixed admissible model calculation, finitely many rational
period coordinates for the actual obstruction. Their vanishing
is a certificate that a specified homological Chow class belongs
to $\cK_K^r(X)$. The normalized tests use
\[
 \eta_j=\frac{1}{[L:K]c_L}\Tr^K_{X_L/X}(b_j),\qquad
 c_L=-\log|\varpi_L|=c_K/e(L/K).
\]
Thus all model families, isotypic sectors and individually certified
classes may be combined inside the same admitted group. For a
number-field class one checks this membership at its finitely many
relevant bad places; the good places are handled by the good-model
construction. The output Green identity still has the complete
actual test space in \eqref{eq:strong-G}.

\subsection{Exact-support cancellation for complete intersections}\label{sub:ci-certificates}

Let $\bar L$ be a line bundle with one fixed piecewise-smooth or
smooth metric allowed by the complete-intersection construction
on a smooth projective $X/K$, and let
$s_0,s_1$ be regular rank-$r$ tuples of sections of $L$, with
$n\geq r\geq3$. Write $Y_i=V(s_i)$ with its scheme-theoretic
cycle multiplicities, and put
\[
 Z=[Y_1]-[Y_0],\quad
 S=|Y_0|^{\an}\cup|Y_1|^{\an},\quad
 U=X^{\an}\setminus|Z|^{\an},\quad W=X^{\an}\setminus S.
\]
The complete-intersection kernel of the tropical companion
\cite{InternalTropical}, for a common allowed metric on $L$, is
\begin{equation}\label{eq:ci-kernel}
 \phi_s=\min_i(-\log\|s_i\|),\quad
 R=c_1(\bar L),\quad A_s=R-D\phi_s,\quad
 G_s=\phi_s\sum_{j=0}^{r-1}A_s^jR^{r-1-j}.
\end{equation}
It satisfies $D[G_s]+\delta_{Y_s}=[R^r]$ against all actual
tests. The raw difference is defined on $W$; cancelled components
lie in $C_U=S\cap U$.

\begin{proposition}[Supplied cancellation certificates]\label{prop:ci-cancellation}
Each of the following hypotheses supplies an explicit actual
Green form for $Z$ on $U$, satisfying \eqref{eq:strong-G}.
\begin{enumerate}[label=\textup{(\roman*)}]
\item On an ordinary analytic neighborhood of $C_U$ in $U$,
$s_1=Ms_0$ exactly, where $M$ is an invertible three-row block
plus the identity in the other rows. Two specified entries $a,d$
in its first column have nonvanishing opens covering $C_U$.
\item On a Zariski neighborhood of $C_U$, the two regular
sequences define the same ideal. Their comparison matrix on the
reduced conormal coefficient ring has a supplied finite ordered
factorization into elementary shears, one-entry units, invertible
two-row blocks, and invertible three-row blocks with the two-pivot
cover in \textup{(i)}. All row labels, coefficient data and inverses
are part of the factorization. At each fixed word position use
the same cap offsets and inner and outer partitions for every
local coefficient lift.
\end{enumerate}
\end{proposition}

\begin{proof}
For (i), the relative three-row face identities in
\cite{InternalTropical} give, for an ordinary neighborhood $N$
of $C_U$ in $U$, the identity on $N\cap W$
\[
 G_{s_1}-G_{s_0}=d'\alpha+d''\beta+E,
\]
where $\alpha,\beta$ are actual on $N\cap W$ and $E$ is
actual on $N$. Their explicit
two-pivot partition uses $v(d)-v(a)$ and includes the differentiated
partition and mixed-derivative terms. Choose the companion
relative-support cutoff $\zeta$ on $U$, with support closed in
$U$ inside this neighborhood and value one near $C_U$.
On $W$ the corrected form
\[
 G_{s_1}-G_{s_0}
 -d'\operatorname{ext}_0(\zeta\alpha)
 -d''\operatorname{ext}_0(\zeta\beta)
\]
equals $E$ near $C_U$, so it glues to an actual form on $U$.
The relative-support Stokes calculation of the companion proves
the all-actual test identity.

For (ii), lift the fixed finite word locally and execute its
factors at the actual intermediate tuples. The final discrepancy
is $1+E_0$ with entries of $E_0$ vanishing on the reduced common
zero scheme. Near $C_U$ its entries have norm less than one, so
this discrepancy preserves the tuple supremum norm and the
complete expression \eqref{eq:ci-kernel}. The coefficient-uniform
germ theorem in \cite{InternalTropical} identifies the complete
primitive germs for different local lifts. They therefore glue
with a locally finite partition; its derivative discrepancy
extends by zero. Apply the same relative-support cutoff argument.
The calculation retains the original nonreduced schemes and their
cycle multiplicities throughout.
\end{proof}

These regular sections have
$[Y_i]=c_r(L^{\oplus r})\cap[X]=c_1(L)^r\cap[X]$ in Chow.
Consequently $[Z]=0$. The constructions make the exact-support
Green representative explicit. For such globally defined cycles,
the Chow-descended height is
\[
 h_X([Z],w)=0\qquad(w\in\cH^{n+1-r}(X)),
\]
even though individual local Green integrals may be nonzero.

\subsection{Comparison with supplied arithmetic models}\label{sub:supplied-models}

Suppose an admitted pair $z,w$ has regular numerically admissible
arithmetic lifts $\widehat\Theta,\widehat\Phi$ on a regular flat
projective model over $\mathcal O_F$, with zero archimedean
curvature in the comparison normalization. The companion
model-comparison theorem \cite{InternalHeight} identifies its
normalized height with
\[
 h_X(z,w)=[F:\Q]^{-1}
          \widehat{\deg}(\widehat\Theta\cdot\widehat\Phi).
\]
If $f:Y\to X_E$ is a degree-$d$ alteration with $Y$ smooth
projective, and the supplied admissible lifts of
$f^*z_E,f^*w_E$ lie on a regular flat projective
$\mathcal O_E$-model of $Y$, with the same zero-curvature
archimedean normalization, the corresponding formula is
\[
 h_X(z,w)=\frac{1}{d[E:\Q]}
          \widehat{\deg}(\widehat\Theta\cdot\widehat\Phi).
\]
The generic degree and the absolute number-field degree are separate
factors. These formulas and the abelian comparison in the next
section describe the same normalized pairing on their common domains.



\section{Height pairings and recursive formulas}\label{sec:heights}

\subsection{Proof of the normalized height theorem}\label{proof:height}

Let $X/F$ satisfy the hypothesis of \thmref{thm:height-main}.
At each finite place $v$, complete a field of definition of
the Chow retraction at a place above $v$. The resulting
extension of $F_v$ is finite. \thmref{thm:main}, applied with
the original absolute value of $F_v$, shows that every geometric
homological class has zero local obstruction. Both groups in
\eqref{eq:height-domain} are therefore in the full domain of
the companion height theorem.

Normalized finite-extension invariance, the Chow identity and
\eqref{eq:height-adjunction} give
\begin{align*}
 h_X^{\an}(z,w)
 &=h_{X_E}^{\an}(z_E,w_E)\\
 &=\sum_jh_{X_E}^{\an}(v_{j*}u_{j*}z_E,w_E)\\
 &=\sum_jh_{A_j}^{\an}(u_{j*}z_E,{}^tv_{j*}w_E).
\end{align*}
The image classes are homologically trivial and have codimensions
$r+a_j$ and $\dim A_j+1-r-a_j$. An image in codimension zero
is zero because its geometric cycle class detects its rank.
Out-of-range images also vanish. Apply
\eqref{eq:abelian-comparison} to each remaining term to obtain
\eqref{eq:height-main}. Every permitted choice computes the
same canonical left side, proving independence. This completes
the proof of \thmref{thm:height-main}.

\subsection{The signed join recursion}

Consider a join $g:W=\Bl_CV\to H$ of degree $d$ over a number
field, with the potential abelian Chow retractions constructed above
for $V,C,W,H$. For geometrically homologically trivial rational
Chow classes $z,w$ on $H$ in complementary codimensions, put
\[
 a_z=q_*g^*z,\quad a_w=q_*g^*w,\qquad
 b_z=p_*i^*g^*z,\quad b_w=p_*i^*g^*w.
\]

\begin{corollary}\label{cor:height-recursion}
The normalized heights satisfy
\begin{equation}\label{eq:height-recursion}
 h_H^{\an}(z,w)=\frac1d
          \left(h_V^{\an}(a_z,a_w)-h_C^{\an}(b_z,b_w)\right).
\end{equation}
The center height is summed over its components after a finite
splitting field. Its two degrees are $r-1$ and $s-1$, where
$r+s=\dim H+1$; zero or impossible homological terms are omitted.
\end{corollary}

\begin{proof}
Adjunction and $g_*g^*=d$ yield
$h_W^{\an}(g^*z,g^*w)=d\,h_H^{\an}(z,w)$. Apply
\eqref{eq:blowup-diagonal} to the first entry on $W$.
The transposes of $q^*$ and $i_*p^*$ act by $q_*$ and
$p_*i^*$, respectively. A second application of adjunction
gives the right side of \eqref{eq:height-recursion}, with
the negative sign from the normal bundle of the exceptional
divisor. The center codimensions sum to $\dim C+1$.
\end{proof}

\subsection{The cyclic-cover scalar}

For the product cover $f:Y_E\to X_E$ of rank $M$ and the
factorization of $Y_E$ used in \eqref{eq:cyclic-diagonal},
the same adjunction gives
\begin{align}
 h_X^{\an}(z,w)
 &=\frac1M h_{Y_E}^{\an}(f^*z_E,f^*w_E)\notag\\
 &=\frac1{M[E:\Q]}\sum_J
   \bigl\langle u_{J*}f^*z_E,{}^tv_{J*}f^*w_E
              \bigr\rangle_{A_J,\mathrm{Kun}}.\label{eq:cyclic-height}
\end{align}
Indeed, $f_*f^*z_E=M z_E$. The global pairing is already
absolutely normalized, while the factor $M^{-1}$ records the
finite flat rank in \eqref{eq:cyclic-algebra}. In particular,
using the maps $(U_J,V_J)$ instead absorbs this factor into
$V_J$ and gives the same formula.

\begin{thebibliography}{99}
\bibitem{ChowGreenPairs} \emph{Chow Obstructions and Prescribed Actual Green Pairs}, companion manuscript, 23 November 2026, Theorem~1.1. \url{https://openmath-atlas.org/projects/strong-green/editions/v0.5.2/papers/chow-obstruction-green-pairs/}.

% Source: papers/pairing-sections/bibliography_additions.tex
\bibitem{DCM} M. A. A. de Cataldo and L. Migliorini,
\emph{The Chow groups and the motive of the Hilbert scheme of points
on a surface}, J. Algebra 251 (2002), 824--848, Theorem 6.2.1 and
Remark 6.2.2. \url{https://arxiv.org/abs/math/0005249}.
\bibitem{Bulles} T.-H. B\"ulles, \emph{Motives of moduli spaces on K3
surfaces and of special cubic fourfolds}, arXiv:1806.08284, Theorem 0.1.
\url{https://arxiv.org/abs/1806.08284}.
\bibitem{VialFibrations} C. Vial, \emph{Algebraic cycles and fibrations},
Documenta Math. 18 (2013), 1521--1553, Theorem 4.2.
\url{https://arxiv.org/abs/1203.2650}.
\bibitem{ELV} H. Esnault, M. Levine and E. Viehweg,
\emph{Chow groups of projective varieties of very small degree},
Duke Math. J. 87 (1997), 29--58, Theorem 4.6.
\url{https://page.mi.fu-berlin.de/esnault/preprints/helene/38-duke97.pdf}.
\bibitem{HI} A. Hirschowitz and J. N. Iyer,
\emph{Hilbert schemes of fat $r$-planes and the triviality of Chow
groups of complete intersections}, final author manuscript,
Theorem 1.5 and Sections 5--6. The exact version is identified in
the source dossier.
\url{https://math.univ-cotedazur.fr/u/ah/div/roitfinalCM.pdf}.
\bibitem{Otwinowska} A. Otwinowska,
\emph{Remarques sur les cycles de petite dimension de certaines
intersections compl\`etes}, C. R. Acad. Sci. Paris S\'er. I Math.
329 (1999), 141--146. Rewritten author version, 13 December 1999,
Corollaire 1.
\url{https://www.imo.universite-paris-saclay.fr/~anna.otwinowska/intcompletes.ps}.
\bibitem{Ancona} G. Ancona, \emph{Standard conjectures for abelian
fourfolds}, Invent. Math. 223 (2021), 149--212.
\bibitem{Ito} T. Ito, \emph{Weight-monodromy conjecture for
$p$-adically uniformized varieties}, Invent. Math. 159 (2005),
607--656. \url{https://arxiv.org/abs/math/0301201}.
\bibitem{JannsenSaito} U. Jannsen and S. Saito,
\emph{Bertini theorems and Lefschetz pencils over discrete valuation
rings, with applications to higher class field theory},
arXiv:0911.1470, Section 3, Lemma 4 and Theorem 5.
\url{https://arxiv.org/abs/0911.1470}.
\bibitem{RX} W. Raskind and X. Xarles,
\emph{On the \'etale cohomology of algebraic varieties with totally degenerate
reduction over $p$-adic fields}, J. Math. Sci. Univ. Tokyo 14
(2007), 261--291. Definition 1(a),(b), Proposition 2 and Corollary 2.
\url{https://www.ms.u-tokyo.ac.jp/journal/pdf/jms140204.pdf}.
\bibitem{InternalEndpoints} Research archive,
\emph{Endpoint cycles, numerical curve lifts, and dominant
obstruction descent}. Exact sources: regular-model curve correction
(R12), all-degree dominant detection and one-cycles (R16).
See the companion source index, group E.
\bibitem{InternalProducts} Research archive,
\emph{Actual deletion models and the residual core for all surface
products}. Exact sources: the binary supported transport, incidence
conventions, and all-factor deletion proof (R20), including the
model-splitting and full local product theorems. Source group P.
\bibitem{InternalModel} Research archive,
\emph{Actual algebraic class spans, incidence defects, and odd-weight
obstructions}. Exact model interfaces AK19, CR19 and OW19;
source group M. The product theorem invokes these interfaces in
their specified actual-class-span form.
\bibitem{InternalFibrations} Research archive,
\emph{Low-Chow fibers and signed quadric retractions}. Exact
quadric-fiber, cubic-fiber, relative-quadric and height-domain
proofs (R32--SC2); source group Q.
\bibitem{InternalModelCriteria} Research archive,
\emph{Primitive residuals, scalar boundary corrections, and model
symmetries}. Exact numerical model proofs (R16); source group M.
\bibitem{InternalOrdinary} Research archive,
\emph{Ordinary quadratic models and exact native normal forms}
(R17), including ON17, the all-characteristic quadric-forest criterion
C217, and the even-dimensional residue-two argument CD217;
arbitrary-thickness extension and global-height
corollaries. Source group O.
\bibitem{InternalUniformization} Research archive,
\emph{Drinfeld uniformization, ramified toroidal strata, and motive
closure} (R20). Source group U.
\bibitem{InternalDense} Research archive,
\emph{Arithmetic congruence pencils and complete-height families}
(R17--R18), with the all-degree and all-finite-extension overlays.
Source group D.
\bibitem{InternalMixedFamilies} Research archive,
\emph{Mixed surface--hyperplane and two-surface arithmetic
families} (R21--R22). Source group D.
\bibitem{InternalPositiveHosts} Research archive,
\emph{Positive fourfold suspensions and supplied projective hosts}
(R22). Source group D.
\bibitem{InternalDegenerate} Research archive,
\emph{Totally degenerate strata and original-source period
annihilation} (R32--SC2), with the component-detection comparison.
Source group M.
\bibitem{InternalDetectors} Research archive,
\emph{Finite actual tests and height-obstruction quotients}
(R17), with exact-support equivalence and native trace weights.
Source group F.
\bibitem{InternalTropical} Research archive,
\emph{Complete-intersection kernels and exact-support
two-pivot and conormal-word cancellation}. Tropical branch,
R9 and R12 exact sources; source group T. The complete source
and dependency audit accompanies the paper.


\bibitem{Mihatsch} A. Mihatsch, \emph{$\delta$-forms on Lubin--Tate spaces},
arXiv:2112.10018v4, 2024. \url{https://arxiv.org/abs/2112.10018}.
\bibitem{Kun} K. K\"unnemann, \emph{Height pairings for algebraic cycles on
abelian varieties}, Ann. Sci. \'Ecole Norm. Sup. (4) 34 (2001), 503--523.
\url{https://www.numdam.org/article/ASENS_2001_4_34_4_503_0.pdf}.
\bibitem{SK} T. Shioda and T. Katsura, \emph{On Fermat varieties},
Tohoku Math. J. (2) 31 (1979), 97--115.
\url{https://doi.org/10.2748/tmj/1178229881}.
\bibitem{Laterveer} R. Laterveer, \emph{Some cubics with finite-dimensional
motive}, arXiv:1708.05671v1, 2017, Theorem 5 and Proposition 6.
\url{https://arxiv.org/abs/1708.05671}.
\bibitem{Vial} C. Vial, \emph{Remarks on motives of abelian type},
Tohoku Math. J. (2) 69 (2017), 195--220; author version of 20 July 2015.
\url{https://www.math.uni-bielefeld.de/~vial/abelian.pdf}.
\bibitem{Scholl} A. J. Scholl, \emph{Classical motives}, in
\emph{Motives}, Proc. Sympos. Pure Math. 55, Part 1 (1994), 163--187.
\url{https://www.dpmms.cam.ac.uk/~ajs1005/preprints/classical_motives.pdf}.
\bibitem{Fulton} W. Fulton, \emph{Intersection Theory}, second edition,
Ergebnisse der Mathematik und ihrer Grenzgebiete 2, Springer, 1998.
\bibitem{InternalGreen} Research archive, \emph{Curved ordinary Green
pairs and the Chow-to-actual Bott--Chern obstruction}, R11 companion
manuscripts, frozen dependencies of R32. Exact statements: PC7;
OC1--OC2; exact-support obstruction criterion. Source paths and
hash/review bindings are supplied in the accompanying source map.
\bibitem{InternalNumerical} Research archive,
\emph{Numerical witness descent}, R12 companion manuscript,
\texttt{r12a:finite-numerical-cover};
\href{https://openmath-atlas.org/projects/strong-green/editions/v0.5.2/source-files/companion/proofs/foundations/R12_numerical_witness_descent.tex}{complete R12 source}.
\bibitem{InternalCovariance} Research archive,
\emph{Chow-kernel covariance and full finite actual trace}, R16 companion
manuscripts: \texttt{r16cov:proper-push}, \texttt{r16cov:correspondences},
\href{https://openmath-atlas.org/projects/strong-green/editions/v0.5.2/source-files/companion/proofs/foundations/R16_proper_covariance_and_rational_domination.tex}{CV16 source};
\texttt{r16ft:full-trace},
\href{https://openmath-atlas.org/projects/strong-green/editions/v0.5.2/source-files/companion/proofs/foundations/R16_full_finite_actual_trace_reconstruction.tex}{FT16 source}.
\bibitem{InternalProjective} Research archive, R16 companion manuscripts:
\emph{All-actual projective-bundle output}, \texttt{r16out:line},
\texttt{r16out:bundle},
\href{https://openmath-atlas.org/projects/strong-green/editions/v0.5.2/source-files/companion/proofs/foundations/R16_all_actual_projective_bundle_output.tex}{OUT16 source};
\emph{All-degree projective cohomology and descent},
\texttt{r16pb:formula}, \texttt{r16pb:blowup},
\href{https://openmath-atlas.org/projects/strong-green/editions/v0.5.2/source-files/companion/proofs/foundations/R16_all_degree_projective_cohomology_and_descent.tex}{PB16 source};
\emph{Smooth-immersion obstruction kernel}, \texttt{r16gi:detection},
\texttt{r16gi:kernel},
\href{https://openmath-atlas.org/projects/strong-green/editions/v0.5.2/source-files/companion/proofs/foundations/R16_smooth_immersion_obstruction_kernel.tex}{GI16 source}.
\bibitem{InternalHeight} Research archive,
\emph{Global analytic heights, correspondence adjunction and abelian
comparison}, R16 companion manuscripts:
\texttt{r16h:curved-slice}, \texttt{r16h:global-height},
\href{https://openmath-atlas.org/projects/strong-green/editions/v0.5.2/source-files/companion/proofs/foundations/R16_obstruction_zero_reciprocity_and_height.tex}{H16 source};
\texttt{r16ha:moving}, \texttt{r16ha:adjunction},
\href{https://openmath-atlas.org/projects/strong-green/editions/v0.5.2/source-files/companion/proofs/foundations/R16_correspondence_adjunction_and_motivic_height.tex}{HA16 source};
\texttt{r16mm:local}, \texttt{r16mm:global}, \texttt{r16mm:abelian},
\href{https://openmath-atlas.org/projects/strong-green/editions/v0.5.2/source-files/companion/proofs/foundations/R16_regular_admissible_model_comparison.tex}{MM16 source}.
\end{thebibliography}
\end{document}
